\section{Random splitting, linking, and the lending colouring}\label{s3:sec}

This section provides the probabilistic tools that every later routing step
uses, together with the stage-1 lending data.
\begin{itemize}
\item The first subsection proves monotonicity facts and
  Lemma~15$^+$. A uniformly random $k$-colouring of the edges of an
  $(\varepsilon',s)$-expander leaves every colour class an
  $(\varepsilon',s/(2k))$-expander: the first parameter is untouched and
  the second loses only the factor $2k$. Lemma~15$^+$ replaces
  \cite[Lemma~15]{BM} (quoted as Cited result~\ref{s1:citLem15}), which is not
  used.
\item The second subsection proves Theorem~16$^*$. It says that a
  robust expander is path connected through a random vertex set, in the
  multiset, for-all form of \cite[Definition~7]{BM}. The requirement on
  the robustness is linear in the multiplicity and polynomial in the
  inverse density. Its proof re-runs the proof of \cite[Theorem~16]{BM}
  (Cited result~\ref{s1:citThm16}) with explicit parameters.
\item The third subsection proves Lemma~HB by one application of Hall's
  theorem.
\item The fourth subsection defines the two-stage lending colouring
  and the JS labels, verifies the numerical conditions they need (the
  COL-JV table), and proves Lemma~COL.
\end{itemize}
Every statement of this section is proved in full here, from the results
stated in Section~\ref{s1:sec} and, in the fourth subsection, from the hierarchy of
Section~\ref{s2:sec}.

\medskip\noindent\textbf{Conventions for this section.} Logarithms are to
base $2$ (Convention~\ref{s1:convGraphs}), and $\ln$ is the natural
logarithm. For $\rho\in[0,1]$, a \emph{$\rho$-random subset} of a finite
set $S$ is a random subset of $S$ that contains each element of $S$
independently with probability $\rho$. The letter $\varepsilon'$ always
denotes a generic expansion parameter in $(0,1]$. The fixed constant
$\eps=2^{-5}$ of Definition~\ref{s1:defConstants} enters this section only
through the values $\varepsilon_Y\in\{2^{-5},2^{-6}\}$ of
Definition~\ref{s2:defAncestors}. For a graph $G$, $U\subseteq V(G)$ and
$F\subseteq E(G)$, the set $\Nbr_{G-F}(U)$ is the outside neighbourhood of
$U$ in the graph $G-F$ (Convention~\ref{s1:convGraphs}). The ball
$\Ball{i}{G-F}{U,V}$ is the set of vertices of $V$ that can be reached
from $U$ by a path of length at most $i$ in $G-F$ whose interior vertices
all lie in $V$ (Cited result~\ref{s1:citNotation}). Path connectivity
(Cited result~\ref{s1:citDef7}) always refers to \emph{multisets} of pairs of
distinct vertices; see the remark on multisets at the end of the first subsection.

%% ----------------------------------------------------------------------
\subsection{Monotonicity and random splitting}

\begin{lemma}[monotonicity and the for-all property]\label{s3:lemMonotone}
\begin{itemize}
\item[(i)] Let $H$ be an $(\varepsilon',s)$-expander and let $H'\supseteq
  H$ be a graph with $V(H')=V(H)$. Then $H'$ is an
  $(\varepsilon',s)$-expander. Moreover, every $(\varepsilon',s)$-expander
  is an $(\varepsilon'',s')$-expander whenever $\varepsilon''\le
  \varepsilon'$ and $s'\le s$.
\item[(ii)] Path connectivity is monotone. Let $G\subseteq G'$ be graphs
  with $V(G)=V(G')$, let $V\subseteq V'\subseteq V(G)$, and suppose
  $\ell\le\ell'$ and $t'\le t$. If $G$ is $(\ell,t)$-path connected
  through $V$, then $G'$ is $(\ell',t')$-path connected through $V'$.
\item[(iii)] Whether a graph $X$ is $(\ell,t)$-path connected through a
  set $V$ depends only on the pair $(X,V)$. On this event, required paths
  exist for \emph{every} multiset $\mathcal P$ of pairs of distinct
  vertices of $X$ in which each vertex lies in at most $t$ pairs. In
  particular, $\mathcal P$ may be chosen after $X$ and $V$ are revealed,
  as an arbitrary function of them and of any further randomness
  (adaptively). If several multisets $\mathcal P_1,\dots,\mathcal P_q$ are
  given and every vertex lies in at most $t$ pairs of the union multiset
  $\mathcal P_1+\dots+\mathcal P_q$, then one application of the property
  to this union joins all their pairs by pairwise edge-disjoint paths
  through $V$ of length at most $\ell$ (joint routing).
\end{itemize}
\deps{\ref{s1:citDef7}, \ref{s1:citDef11}}
\end{lemma}

\begin{proof}
(i) Put $n:=|V(H)|$. Let $U\subseteq V(H')$ with $1\le|U|\le 2n/3$, and
let $F\subseteq E(H')$ with $|F|\le s|U|$. Then $F\cap E(H)$ is a set of
at most $s|U|$ edges of $H$, and $H-(F\cap E(H))\subseteq H'-F$. Hence
\[
|\Nbr_{H'-F}(U)|\ \ge\ |\Nbr_{H-(F\cap E(H))}(U)|\ \ge\ \varepsilon'|U|/\log^2 n
\]
by Definition~11 of \cite{BM} (Cited result~\ref{s1:citDef11}) applied to $H$.
Both graphs have $n$ vertices, so the logarithm in the bound is the same.
The second claim holds because the condition of Cited result~\ref{s1:citDef11}
becomes weaker when $\varepsilon'$ or $s$ decreases.

(ii) Let $\mathcal P$ be a multiset of pairs of distinct vertices of
$G'$ in which every vertex lies in at most $t'\le t$ pairs. The
hypothesis on $G$ applies to $\mathcal P$, since $V(G)=V(G')$. It gives
pairwise edge-disjoint paths of $G$, one for each pair, each of length at
most $\ell\le\ell'$ and with all interior vertices in $V\subseteq V'$.
These are paths of $G'$ with the required properties.

(iii) By Cited result~\ref{s1:citDef7}, the property says that for every
admissible multiset the paths exist. It mentions no other data, so it is
a property of $(X,V)$ alone, and on this event it applies to every
admissible multiset, however that multiset was chosen. The union
$\mathcal P_1+\dots+\mathcal P_q$ is itself an admissible multiset, which
gives the last sentence.
\end{proof}

\begin{lemma}[Lemma 15$^+$: random splitting without loss]\label{s3:lemL15p}
Let $X$ be an $N$-vertex $(\varepsilon',s)$-expander with
$0<\varepsilon'\le1$, put $L:=\log N$, and let $k\ge1$ be an integer with
$s\ge 40kL$. Give every edge of $X$ a colour from $[k]$, independently
and uniformly at random, and for $i\in[k]$ let $X_i$ be the graph with
vertex set $V(X)$ whose edges are the edges of colour $i$. Then for each
$i\in[k]$, $X_i$ is an $(\varepsilon',s/(2k))$-expander with probability
at least $1-2N^{-5}$. Hence all $k$ classes are
$(\varepsilon',s/(2k))$-expanders with probability at least $1-2kN^{-5}$.
\deps{\ref{s1:citDef11}, \ref{s1:citChernoff}}
\end{lemma}

\begin{proof}
If $N=1$, then no set $U$ satisfies $1\le|U|\le 2N/3$, so every graph on
one vertex is an expander for all parameters, and there is nothing to
prove. Assume $N\ge2$, so $L\ge1$. Fix $i\in[k]$ and write
$H:=X_i$.

Fix $U\subseteq V(X)$ with $1\le u:=|U|\le 2N/3$. Put
$m:=\lceil\varepsilon' u/L^2\rceil$ and
$\mathcal N:=\Nbr_X(U)$. Since $\varepsilon'\le1\le L^2$, we have
$1\le m\le u$.

\smallskip\noindent\emph{Step 1 (every small deletion leaves many edges).}
We claim that
\begin{equation}\label{s3:eqL15step1}
e_X(U,\mathcal N\setminus\mathcal T)>su\qquad\text{for every }\mathcal T\subseteq\mathcal N\text{ with }|\mathcal T|\le m-1 .
\end{equation}
Suppose instead that $e_X(U,\mathcal N\setminus\mathcal T)\le su$ for some
such $\mathcal T$, and let $F$ be the set of edges of $X$ between $U$ and
$\mathcal N\setminus\mathcal T$. Then $|F|\le su$, and every neighbour of
$U$ in $X-F$ outside $U$ lies in $\mathcal T$. Hence
$|\Nbr_{X-F}(U)|\le m-1<\varepsilon' u/L^2$, which contradicts the
expansion of $X$ (Cited result~\ref{s1:citDef11}).

\smallskip\noindent\emph{Step 2 (what a failure looks like).} Suppose
that $U$ violates the $(\varepsilon',s/(2k))$-condition in $H$. Then
there is $F\subseteq E(H)$ with $|F|\le su/(2k)$ and
$|\Nbr_{H-F}(U)|<\varepsilon'u/L^2$. Since $|\Nbr_{H-F}(U)|$ is an
integer, it is at most $m-1$. Put $\mathcal T:=\Nbr_{H-F}(U)$. Since
$H\subseteq X$, we have $\mathcal T\subseteq\mathcal N$ and
$|\mathcal T|\le m-1$, and every edge of $H$ between $U$ and
$\mathcal N\setminus\mathcal T$ lies in $F$. Hence
$e_H(U,\mathcal N\setminus\mathcal T)\le su/(2k)$.

\smallskip\noindent\emph{Step 3 (one pair $(U,\mathcal T)$).} Fix
$\mathcal T\subseteq\mathcal N$ with $|\mathcal T|\le m-1$. Since $U$ and
$\mathcal N\setminus\mathcal T$ are disjoint, the edges of $X$ between
them are $e:=e_X(U,\mathcal N\setminus\mathcal T)$ distinct edges, and
$e>su$ by \eqref{s3:eqL15step1}. Each of them lies in $H$ independently
with probability $1/k$. So
$Z:=e_H(U,\mathcal N\setminus\mathcal T)\sim\Bin(e,1/k)$, with mean
$\mu_Z:=e/k>su/k$. The event $\{Z\le su/(2k)\}$ is contained in
$\{Z<\mu_Z/2\}$. By the lower Chernoff bound (Cited result~\ref{s1:citChernoff})
with $\delta=1/2$,
\[
\Prob\big(Z\le su/(2k)\big)\le e^{-\mu_Z/8}\le e^{-su/(8k)}\le e^{-5uL},
\]
where the last step uses $s\ge 40kL$.

\smallskip\noindent\emph{Step 4 (union bound).} There are at most $N^u$
sets $U$ of size $u$. The number of sets $\mathcal T$ with
$|\mathcal T|\le m-1$ is at most $\sum_{j<m}N^j\le N^m\le N^u$. Since
$e^{-5uL}=N^{-5u\log e}\le N^{-7.2u}$, the probability that $H$ is not an
$(\varepsilon',s/(2k))$-expander is at most
\[
\sum_{u\ge1}N^{2u}\cdot N^{-7.2u}\ \le\ \frac{N^{-5.2}}{1-N^{-5.2}}\ \le\ 2N^{-5}.
\]
The final statement follows by a union bound over the $k$ colours.
Neither $\varepsilon'$ nor the logarithm $L$ changes, because $X_i$ has
the same vertex set as $X$.
\end{proof}

\begin{remark}[multisets; \RTb{I14}]\label{s3:remMultiset}
In \cite[Definition~7]{BM} (Cited result~\ref{s1:citDef7}), the collection
$\mathcal P$ is a \emph{multiset} of pairs of distinct vertices.
\begin{itemize}
\item The same pair may occur several times, and each occurrence needs a
  path of its own.
\item The endpoints are arbitrary vertices of the graph; they need not
  lie in $V$.
\item The proof of \cite[Lemma~9]{BM} (Cited result~\ref{s1:citLem9}) indexes
  the pairs by $i\in[r]$, so it applies verbatim to multisets.
\item Its counting step also holds when pairs are counted with
  multiplicity. The step says that a maximal subfamily $I'\subseteq I$ of
  pairwise vertex-disjoint pairs satisfies $2t|I'|\ge|I|$. It holds
  because every pair of $I$ meets one of the $2|I'|$ vertices covered by
  $I'$, and each vertex lies in at most $t$ pairs, counted with
  multiplicity.
\end{itemize}
All path-connectivity statements of this manuscript are meant in this
multiset sense, and the proofs of this subsection and the next treat the
pairs as an indexed family. Repeated pairs genuinely occur later: in the
routing inside the vortex devices of Section~4, and in the joint routing
of the JS-LC systems in Section~6.
\deps{\ref{s1:citDef7}, \ref{s1:citLem9}}
\fixes{\RTb{I14}}
\end{remark}

%% ----------------------------------------------------------------------
\subsection{Linking through random sets: Theorem 16\texorpdfstring{$^*$}{*}}

This subsection proves Theorem~16$^*$, stated at its end. The route is
that of \cite[Section~4]{BM}:
\begin{enumerate}
\item well-expanding sets grow into a random set by sprinkling
  (Lemma~17$^*$ below, modelled on \cite[Lemma~17]{BM});
\item every set contains a well-expanding subset (Proposition~18$^*$
  and Lemma~19$^*$ below, modelled on \cite[Proposition~18 and
  Lemma~19]{BM});
\item ball expansion implies multiset linking (Lemma~9$_\rho$ below,
  modelled on \cite[Lemma~9]{BM});
\item a random edge-splitting (Lemma~\ref{s3:lemL15p}) assembles these.
\end{enumerate}
Every parameter is explicit, and the density $\rho$ and the multiplicity
$t$ are free.

\medskip\noindent\textbf{Local parameters.} Let $n\ge2$ be the number of
vertices of the graph at hand, $L:=\log n$, $\varepsilon'\in[2^{-7},1]$
its expansion parameter, $\rho\in(0,1]$ a density and $t\ge1$ a
multiplicity. The following quantities are local to this subsection. All
carry a star subscript; the letters $\ell$, $g$, $p$, $q$, $d$,
$\lambda$, $\Delta$, $\sigma$, $\theta$, $\mu$, $\bar s$, $K$, $M$, $t_0$
without a star keep the meaning they have elsewhere. Throughout,
$\ell_*$, $d_*$, $\lambda_*$, $\Delta_*$, $M_*$ and $K_*$ are integers.
\begin{equation}\label{s3:eqStar}
\begin{aligned}
&\ell_*:=\lfloor 2^{10}L^3\rfloor,\qquad g_*:=\frac{\varepsilon'}{8L^2},\qquad q_*:=0.9\rho,\qquad
 p_*\in(0,1]\ \text{ given by }\ (1-p_*)^{\ell_*-1}=\frac{1-\rho}{1-0.9\rho},\\
&d_*:=\lceil 1/p_*\rceil,\qquad \lambda_*:=\lceil 6/p_*\rceil,\qquad
 \Delta_*:=\lambda_*+\lceil d_*\lambda_*/3\rceil,\qquad \sigma_*:=\frac{24L^2}{\varepsilon'},\\
&\theta_*:=\frac{2^{19}\ell_*^2L^3}{\varepsilon'\rho^2},\qquad
 \mu_*:=\frac{\varepsilon'}{6\theta_*L^2},\qquad
 \bar s_*:=\frac{2^{28}tL^8}{\rho},\qquad
 M_*:=\lceil 2.1/\rho\rceil,\qquad
 K_*:=\Big\lceil\frac{6\,\bar s_*\,\theta_*L^2}{\varepsilon'}\Big\rceil .
\end{aligned}
\end{equation}
The integer $t_{0*}$ is defined inside the proof of Lemma~9$_\rho$
below. We record six elementary facts, each with its
proof. Write $p:=p_*$.
\begin{itemize}
\item[(S1)]\namedlabel{s3:eqS1}{(S1)} $2^{10}L^3-1<\ell_*\le 2^{10}L^3$ and $\ell_*\ge 2^{10}$.
  \emph{Proof:} $L\ge1$ because $n\ge2$.
\item[(S2)]\namedlabel{s3:eqS2}{(S2)} The number $p_*$ exists, is unique, and satisfies $p_*\ge 0.1\rho/\ell_*$, so
  $p_*^{-2}\le 100\ell_*^2\rho^{-2}$. If $V_1,\dots,V_{\ell_*}$ are independent random subsets of a
  set $S$, where $V_i$ is $p_*$-random for $i<\ell_*$ and $V_{\ell_*}$ is $q_*$-random, then
  $V_1\cup\dots\cup V_{\ell_*}$ is a $\rho$-random subset of $S$.
  \emph{Proof:} since $\ell_*\ge2$, the map $x\mapsto(1-x)^{\ell_*-1}$ is a continuous decreasing
  bijection of $[0,1]$ onto $[0,1]$, and the value $(1-\rho)/(1-0.9\rho)$ lies in $[0,1)$. Hence
  $p_*\in(0,1]$ exists and is unique. By Bernoulli's inequality,
  $1-(\ell_*-1)p\le(1-p)^{\ell_*-1}=1-0.1\rho/(1-0.9\rho)$. Hence $(\ell_*-1)p\ge 0.1\rho$, which
  gives $p\ge 0.1\rho/(\ell_*-1)\ge 0.1\rho/\ell_*$. For the last claim, an element avoids the union
  with probability $(1-p)^{\ell_*-1}(1-0.9\rho)=1-\rho$, independently over elements.
\item[(S3)]\namedlabel{s3:eqS3}{(S3)} $\Delta_*\le 2p^{-2}+8.34p^{-1}+2.34$. Hence
  $\Delta_*\le 3p^{-2}$ if $p\le 0.11$, and $\Delta_*\le 13p^{-2}$ for every $p\in(0,1]$.
  \emph{Proof:} $d_*\le p^{-1}+1$ and $\lambda_*\le 6p^{-1}+1$, so
  $d_*\lambda_*\le 6p^{-2}+7p^{-1}+1$. Then
  $\Delta_*\le \lambda_*+d_*\lambda_*/3+1\le(6p^{-1}+1)+(2p^{-2}+\tfrac73p^{-1}+\tfrac13)+1
  =2p^{-2}+\tfrac{25}{3}p^{-1}+\tfrac73$.
  If $p\le0.11$, then $8.34p^{-1}\le 0.92p^{-2}$ and $2.34\le 0.03p^{-2}$. If $p\le1$, then
  $8.34p^{-1}+2.34\le 10.68p^{-2}$.
\item[(S4)]\namedlabel{s3:eqS4}{(S4)} $8d_*\lambda_*\le 112p^{-2}\le\theta_*$ and
  $\theta_*\le 2^{39}L^9/(\varepsilon'\rho^2)\le 2^{46}L^9\rho^{-2}$.
  \emph{Proof:} $8d_*\lambda_*\le 8(6+7+1)p^{-2}$ since $p\le1$. By~\ref{s3:eqS2},
  $112p^{-2}\le 11200\,\ell_*^2\rho^{-2}\le 2^{19}\ell_*^2L^3/(\varepsilon'\rho^2)$. The upper bound
  uses $\ell_*^2\le 2^{20}L^6$ and $\varepsilon'\ge2^{-7}$.
\item[(S5)]\namedlabel{s3:eqS5}{(S5)} $K_*\ge\bar s_*/\mu_*$;
  $K_*\le(6\cdot2^{81}+1)\,tL^{19}\rho^{-3}\le 2^{83.6}\,tL^{19}\rho^{-3}$;
  $2K_*(\theta_*+1)\le 2^{130.6}\,tL^{28}\rho^{-5}$; and $40K_*L\le 2K_*(\theta_*+1)$.
  \emph{Proof:} $\bar s_*/\mu_*=6\bar s_*\theta_*L^2/\varepsilon'\le K_*$. By~\ref{s3:eqS4},
  $6\bar s_*\theta_*L^2/\varepsilon'\le 6\cdot2^{28}tL^8\rho^{-1}\cdot 2^{39}L^9\rho^{-2}\cdot L^2\cdot
  2^{14}=6\cdot2^{81}tL^{19}\rho^{-3}$, since $\varepsilon'^{-2}\le2^{14}$. The rounding adds at most
  $1\le tL^{19}\rho^{-3}$. Likewise $\theta_*+1\le(2^{46}+1)L^9\rho^{-2}$, and
  $\log\big(2(6\cdot2^{81}+1)(2^{46}+1)\big)<130.6$. Finally
  $\theta_*\ge 2^{19}L^3\ge 20L$.
\item[(S6)]\namedlabel{s3:eqS6}{(S6)} $\sigma_*\ge24$, $\lambda_*p\ge6$ and $1\le d_*p\le 1+p\le2$.
  \emph{Proof:} immediate from~\eqref{s3:eqStar}, since $\varepsilon'\le1\le L$.
\end{itemize}
The values $\sigma_*,\lambda_*,\Delta_*,d_*$ of~\eqref{s3:eqStar} satisfy the parameter
hypotheses of Proposition~13$^*$ below: $\sigma_*=24L^2/\varepsilon'$ holds with equality, and
$\Delta_*-\lambda_*=\lceil d_*\lambda_*/3\rceil\ge d_*\lambda_*/3=8d_*\lambda_*L^2/(\varepsilon'\sigma_*)$.
The proposition itself is stated for arbitrary numbers with these properties.

\begin{proposition}[Proposition 13$^*$: stars or bipartite]\label{s3:propP13s}
Let $G$ be an $n$-vertex $(\varepsilon',s_1)$-expander with $n\ge2$ and
$2^{-7}\le\varepsilon'\le1$, and put $L:=\log n$. Let
$W\subseteq V(G)$ with $1\le|W|\le 2n/3$, and let $F\subseteq E(G)$ with
$|F|\le s_1|W|/4$. Let $\sigma_*>0$ be real and let
$\lambda_*,d_*\ge1$ and $\Delta_*\ge\lambda_*$ be integers such that
\[
\sigma_*\ge\frac{24L^2}{\varepsilon'},\qquad
\Delta_*-\lambda_*\ge\frac{8d_*\lambda_*L^2}{\varepsilon'\sigma_*},\qquad
s_1\ge 8d_*\lambda_* .
\]
Then $G-F$ contains one of the following:
\begin{itemize}
\item[(a)] at least $|W|/\sigma_*$ vertex-disjoint stars, each with
  exactly $\lambda_*$ leaves, its centre in $W$ and all its leaves in
  $V(G)\setminus W$; or
\item[(b)] a bipartite subgraph $H\subseteq G-F$ with sides $W$ and
  $X\subseteq V(G)\setminus W$, such that $|X|\ge\varepsilon'|W|/(2L^2)$,
  every vertex of $X$ has degree exactly $d_*$ in $H$, and every vertex
  of $W$ has degree at most $\Delta_*$ in $H$.
\end{itemize}
\deps{\ref{s1:citProp12}, \ref{s1:citProp13}, \ref{s1:citDef11}}
\end{proposition}

\begin{proof}
We follow the proof of \cite[Proposition~13]{BM}
(Cited result~\ref{s1:citProp13}) with general parameters. For $U\subseteq
V(G)$ and an integer $d\ge1$, write $\Nbr_{G-F,d}(U)$ for the set of
vertices outside $U$ that have at least $d$ neighbours in $U$ in the
graph $G-F$. We use \cite[Proposition~12]{BM} (Cited result~\ref{s1:citProp12}):
if $|U|\le 2n/3$, $|F|\le s_1|U|/2$ and $0<d\le s_1$, then
$|\Nbr_{G-F}(U)|\ge s_1|U|/(2d)$ or
$|\Nbr_{G-F,d}(U)|\ge\varepsilon'|U|/L^2$.

Assume that (a) fails. Take a maximal collection of vertex-disjoint stars
in $G-F$, each with exactly $\lambda_*$ leaves, its centre in $W$ and its
leaves in $V(G)\setminus W$. Let $\mathrm{Cen}\subseteq W$ be the set of
centres and $\mathrm{Lvs}\subseteq V(G)\setminus W$ the set of leaves.
Since (a) fails, $|\mathrm{Cen}|<|W|/\sigma_*$, and
$|\mathrm{Lvs}|=\lambda_*|\mathrm{Cen}|<\lambda_*|W|/\sigma_*$.
By maximality, every $w\in W\setminus\mathrm{Cen}$ has at most
$\lambda_*-1$ neighbours in $G-F$ in $V(G)\setminus(W\cup\mathrm{Lvs})$.
(Otherwise a new star with centre $w$ and $\lambda_*$ such leaves could be
added. It is disjoint from the others, because $w\notin\mathrm{Cen}$,
$w\notin\mathrm{Lvs}$ and its leaves avoid $W\cup\mathrm{Lvs}$.) Every
vertex of $\Nbr_{G-F}(W\setminus\mathrm{Cen})$ lies in $\mathrm{Cen}$, in
$\mathrm{Lvs}$, or in $V(G)\setminus(W\cup\mathrm{Lvs})$. Hence
\begin{equation}\label{s3:eqP13a}
|\Nbr_{G-F}(W\setminus\mathrm{Cen})|
<\frac{(1+\lambda_*)|W|}{\sigma_*}+(\lambda_*-1)|W|
\le\Big(\frac{1+\lambda_*}{24}+\lambda_*-1\Big)|W| ,
\end{equation}
where we used $\sigma_*\ge 24L^2/\varepsilon'\ge24$.

\smallskip\noindent\emph{The greedy bipartite graph.} Label the vertices
of $V(G)\setminus W$ as $v_1,\dots,v_{r'}$ and let $H_0$ be the graph on
$W$ with no edges. For $i=1,\dots,r'$, check whether $v_i$ has $d_*$
neighbours in $G-F$ inside $W$, each of degree less than $\Delta_*$ in
$H_{i-1}$. If so, let $H_i$ be $H_{i-1}$ together with $v_i$ and these
$d_*$ edges, and put $v_i$ into $X$. Otherwise put $H_i:=H_{i-1}$.
Finally let $H:=H_{r'}$.
\begin{itemize}
\item $H\subseteq G-F$ is bipartite with sides $W$ and $X$.
\item Every vertex of $X$ has degree exactly $d_*$ in $H$.
\item Every vertex of $W$ has degree at most $\Delta_*$ in $H$: a degree
  increases, by one, only when it is below the integer $\Delta_*$.
\end{itemize}
It remains to show $|X|\ge\varepsilon'|W|/(2L^2)$.

\smallskip\noindent\emph{Saturated vertices.} Let
$\mathrm{Sat}:=\{w\in W\setminus\mathrm{Cen}:\deg_H(w)=\Delta_*\}$. A
vertex $w\in\mathrm{Sat}$ has at most $\lambda_*-1$ neighbours of $G-F$
in $V(G)\setminus(W\cup\mathrm{Lvs})$. Hence it has at most
$\lambda_*-1$ $H$-neighbours in $X\setminus\mathrm{Lvs}$, and so more
than $\Delta_*-\lambda_*$ $H$-neighbours in $X\cap\mathrm{Lvs}$. Every
vertex of $X$ has exactly $d_*$ $H$-neighbours, so
\[
(\Delta_*-\lambda_*)|\mathrm{Sat}|\le e_H(\mathrm{Sat},X\cap\mathrm{Lvs})\le d_*|\mathrm{Lvs}|<\frac{d_*\lambda_*|W|}{\sigma_*}.
\]
The hypothesis on $\Delta_*-\lambda_*$ forces $\Delta_*>\lambda_*$, and it
gives $|\mathrm{Sat}|<\varepsilon'|W|/(8L^2)$. Put
$\mathrm{Bl}:=\mathrm{Cen}\cup\mathrm{Sat}$. Then
\begin{equation}\label{s3:eqP13b}
|\mathrm{Bl}|<\frac{|W|}{\sigma_*}+\frac{\varepsilon'|W|}{8L^2}\le\frac{\varepsilon'|W|}{24L^2}+\frac{\varepsilon'|W|}{8L^2}=\frac{\varepsilon'|W|}{6L^2}\le\frac{|W|}{6},
\end{equation}
so $|W\setminus\mathrm{Bl}|\ge 5|W|/6$.

\smallskip\noindent\emph{Applying Proposition~12.} The set
$U:=W\setminus\mathrm{Bl}$ has $1\le|U|\le 2n/3$, and
$|F|\le s_1|W|/4\le s_1|U|/2$. Also $1\le d_*\le s_1$, since
$s_1\ge 8d_*\lambda_*$. By Cited result~\ref{s1:citProp12} one of two cases
holds.

\emph{Case 1: $|\Nbr_{G-F}(U)|\ge s_1|U|/(2d_*)$.} Then
$|\Nbr_{G-F}(U)|\ge s_1|W|/(3d_*)\ge 8\lambda_*|W|/3$. Let
$x\in\Nbr_{G-F}(U)$. Either $x\notin W\setminus\mathrm{Cen}$, and then
$x\in\Nbr_{G-F}(W\setminus\mathrm{Cen})$; or $x\in W\setminus\mathrm{Cen}$
with $x\notin U$, and then $x\in\mathrm{Sat}\subseteq W$. Hence
$|\Nbr_{G-F}(W\setminus\mathrm{Cen})|\ge(8\lambda_*/3-1)|W|$. Comparing
with~\eqref{s3:eqP13a} gives $8\lambda_*/3<(1+\lambda_*)/24+\lambda_*$,
that is, $39\lambda_*<1$. This contradicts $\lambda_*\ge1$.

\emph{Case 2: $|\Nbr_{G-F,d_*}(U)|\ge\varepsilon'|U|/L^2$.} Let
$v\in\Nbr_{G-F,d_*}(U)\setminus\mathrm{Bl}$. Then $v\notin W$, since
$v\notin U$ and $v\notin\mathrm{Bl}$, and $v$ has at least $d_*$
neighbours of $G-F$ in $U$. Vertices of $U$ lie in
$W\setminus\mathrm{Cen}$ and not in $\mathrm{Sat}$, so they have
$H$-degree below $\Delta_*$, and their degree in $H_{i-1}$ is at most
that. So when $v=v_i$ was processed, a star was available, and $v\in X$.
Therefore, using~\eqref{s3:eqP13b},
\[
|X|\ge|\Nbr_{G-F,d_*}(U)|-|\mathrm{Bl}|\ge\frac{\varepsilon'}{L^2}\cdot\frac{5|W|}{6}-\frac{\varepsilon'|W|}{6L^2}=\frac{2\varepsilon'|W|}{3L^2}\ge\frac{\varepsilon'|W|}{2L^2}.
\]
So (b) holds.
\end{proof}

\begin{lemma}[Lemma 17$^*$: sprinkling at density $\rho$]\label{s3:lemL17s}
Let $G$ be an $n$-vertex $(\varepsilon',s_1)$-expander with $n\ge2$ and
$2^{-7}\le\varepsilon'\le1$, let $\rho\in(0,1]$, take the parameters
of~\eqref{s3:eqStar} with $L:=\log n$, and assume $s_1\ge 8d_*\lambda_*$.
Let $V_1,\dots,V_{\ell_*}$ be independent random subsets of $V(G)$,
where $V_i$ is $p_*$-random for $i<\ell_*$ and $V_{\ell_*}$ is
$q_*$-random, and put $V:=V_1\cup\dots\cup V_{\ell_*}$. By~\ref{s3:eqS2},
$V$ is a $\rho$-random subset of $V(G)$. Call $U\subseteq V(G)$
\emph{well-expanding} if $|\Nbr_G(U)|\ge\theta_*|U|$. Then for every
well-expanding $U$ and every $F\subseteq E(G)$ with $|F|\le|U|$,
\[
\Prob\Big(|\Ball{\ell_*}{G-F}{U,V}|\le|V|/2\Big)\le\exp(-7|U|L).
\]
Since this event depends only on $V$, the same bound holds for every
$\rho$-random subset $V$ of $V(G)$.
\deps{\ref{s3:propP13s}, \ref{s1:citLem17}, \ref{s1:lemBBD}, \ref{s1:citChernoff}, \ref{s1:citChernoffGen}}
\end{lemma}

\begin{proof}
The structure is that of the proof of \cite[Lemma~17]{BM}
(Cited result~\ref{s1:citLem17}). The batch probabilities are adapted to the
density $\rho$, and the concentration step in case~(b) uses the Bernstein
inequality for bounded differences (Lemma~\ref{s1:lemBBD}).

If $U=\emptyset$ the bound is $1$, so let $u:=|U|\ge1$.

\smallskip\noindent\emph{Step 0: $n$ and $\rho n$ are large.} Since $U$
is well-expanding, $\theta_*u\le|\Nbr_G(U)|<n$. Hence
$n>\theta_*\ge 2^{19}\ell_*^2L^3/\rho^2$, and so
$\rho n\ge\rho^2n>2^{19}\ell_*^2L^3\ge 2^{39}$ by~\ref{s3:eqS1}. Also
$7uL<7nL/\theta_*\le 7\rho^2n/(2^{19}\ell_*^2L^2)\le 2^{-36}\rho n$.

\smallskip\noindent\emph{Step 1: the reached sets.} For
$0\le i\le\ell_*$, let $\mathcal B_i$ be the set of vertices $v\in V(G)$
that can be reached from $U$ by a path in $G-F$ of length at most $i$
whose interior vertices (if any) lie in $V_1\cup\dots\cup V_{i-1}$. The
vertices of $\mathcal B_i$ themselves need not lie in any $V_j$. Then
$\mathcal B_0=U$, $\mathcal B_1=U\cup\Nbr_{G-F}(U)$ and
$\mathcal B_0\subseteq\mathcal B_1\subseteq\dots\subseteq\mathcal
B_{\ell_*}$. The set $\mathcal B_i$ is determined by $U$, $F$ and
$V_1,\dots,V_{i-1}$, so it is independent of $V_i,\dots,V_{\ell_*}$.
\begin{itemize}
\item[(O1)] Let $v\in\Nbr_{G-F}(\mathcal B_i)$ have a neighbour $w$ in
  $G-F$ with $w\in\mathcal B_i\cap V_i$. Then
  $v\in\mathcal B_{i+1}$. Indeed, take a path from $U$ to $w$ as in the
  definition of $\mathcal B_i$. The vertex $v$ does not lie on it, since
  otherwise an initial segment would put $v$ into $\mathcal B_i$.
  Appending the edge $wv$ gives a path of length at most $i+1$ whose
  interior lies in $V_1\cup\dots\cup V_i$.
\item[(O2)] $\mathcal B_{\ell_*}\cap V_{\ell_*}\subseteq\Ball{\ell_*}{G-F}{U,V}$,
  since the interiors lie in $V$ and the end lies in $V_{\ell_*}\subseteq V$.
\item[(O3)] $|\mathcal B_1|\ge\theta_*u$. Deleting one edge removes at
  most one vertex from an outside neighbourhood, so
  $|\Nbr_{G-F}(U)|\ge|\Nbr_G(U)|-|F|\ge(\theta_*-1)u$.
\end{itemize}

\smallskip\noindent\emph{Step 2: growth into a fixed set.} Let
$1\le i\le\ell_*-1$, and let $W\subseteq V(G)$ be a fixed set with
$\mathcal B_1\subseteq W$ and $|W|\le 2n/3$. Put
\[
A_i(W):=\{v\in\Nbr_{G-F}(W):\ v\text{ has a neighbour in }G-F\text{ that lies in }W\cap V_i\}.
\]
We claim that
\begin{equation}\label{s3:eqL17claim}
\Prob\big(|A_i(W)|<g_*|W|\big)\le e^{-18uL}.
\end{equation}
By (O3), $|W|\ge\theta_*u$. Also $|F|\le u\le|W|\le s_1|W|/4$, because
$s_1\ge8d_*\lambda_*\ge48$. So Proposition~\ref{s3:propP13s} applies to
$W$ and $F$ with the values~\eqref{s3:eqStar}, and gives one of two
cases.

\emph{Case (a): stars.} Let $\mathrm{Cen}$ be the set of centres of at
least $|W|/\sigma_*$ vertex-disjoint stars as in (a). Every leaf of a
star whose centre lies in $V_i$ belongs to $A_i(W)$, and distinct stars
have distinct leaves. Now $|\mathrm{Cen}\cap V_i|\sim\Bin(|\mathrm{Cen}|,p_*)$,
so the Chernoff bound (Cited result~\ref{s1:citChernoff}, $\delta=1/2$) gives
$\Prob(|\mathrm{Cen}\cap V_i|<p_*|\mathrm{Cen}|/2)\le\exp(-p_*|\mathrm{Cen}|/8)$.
Outside this event, by~\ref{s3:eqS6},
\[
|A_i(W)|\ge\lambda_*p_*|\mathrm{Cen}|/2\ge3|\mathrm{Cen}|\ge3|W|/\sigma_*=g_*|W|.
\]
For the exponent, \ref{s3:eqS2} and $|W|\ge\theta_*u$ give
\[
\frac{p_*|\mathrm{Cen}|}{8}\ge\frac{p_*\varepsilon'|W|}{192L^2}\ge\frac{0.1\rho}{\ell_*}\cdot\frac{2^{19}\ell_*^2L^3u}{192\,\rho^2L^2}\ge 273\,\ell_*uL\ge 18uL.
\]

\emph{Case (b): a bipartite graph.} Let $H\subseteq G-F$ be bipartite
with sides $W$ and $X$ as in (b). Every $x\in X$ with an $H$-neighbour in
$V_i$ lies in $A_i(W)$, because $x\notin W$ and $H\subseteq G-F$. Let
$Z:=|\{x\in X:\Nbr_H(x)\cap V_i\ne\emptyset\}|$. For $x\in X$,
by~\ref{s3:eqS6},
$\Prob(\Nbr_H(x)\cap V_i=\emptyset)=(1-p_*)^{d_*}\le e^{-p_*d_*}\le e^{-1}$,
so $\Ex Z\ge(1-e^{-1})|X|\ge0.63|X|$.

The count $Z$ is a function of the independent indicators
$\xi_w:=\ind{w\in V_i}$, $w\in W$, each Bernoulli with parameter $p_*$.
Order $W$ as $w_1,\dots,w_{|W|}$. For $y=(y_1,\dots,y_{|W|})\in\{0,1\}^{|W|}$
let $\Psi(y)$ be the number of $x\in X$ such that $y_k=1$ for some $k$ with
$w_k\in\Nbr_H(x)$. Since $H$ is bipartite with sides $W$ and $X$, we have
$\Nbr_H(x)\subseteq W$ for $x\in X$, so
$Z=\Psi(\xi_{w_1},\dots,\xi_{w_{|W|}})$. If $y,y'\in\{0,1\}^{|W|}$ differ
only in the $k$-th coordinate, then only the vertices $x\in\Nbr_H(w_k)$ can
change their contribution to $\Psi$, each by at most $1$, so
$|\Psi(y)-\Psi(y')|\le\deg_H(w_k)\le\Delta_*$ by
Proposition~\ref{s3:propP13s}(b). Using $1-p_*\le1$,
$\deg_H(w_k)\le\Delta_*$, $\sum_{w\in W}\deg_H(w)=e(H)=d_*|X|$ (every
$x\in X$ has exactly $d_*$ neighbours in $H$, all in $W$) and~\ref{s3:eqS6},
\[
\sum_k p_*(1-p_*)\deg_H(w_k)^2\le p_*\Delta_*\sum_{w\in W}\deg_H(w)=p_*d_*\Delta_*|X|\le 2\Delta_*|X|=:\beta ,
\]
and $\beta>0$, since $|X|\ge\varepsilon'|W|/(2L^2)>0$ (as $|W|\ge\theta_*u>0$)
and $\Delta_*\ge\lambda_*\ge1$.
Apply the second bound of Lemma~\ref{s1:lemBBD} with $m:=|W|$,
$I_k:=\xi_{w_k}$, $p_k:=p_*$, $c_k:=\deg_H(w_k)$, with $\Delta_*$ in the role
of the bound $b$ of that lemma, this $\beta$ and $a:=0.3|X|$. Since
$2(\beta+\Delta_*a/3)=4\Delta_*|X|+0.2\Delta_*|X|$ and
$0.09/4.2>1/47$, it gives
\[
\Prob\big(Z\le\Ex Z-0.3|X|\big)\le\exp\Big(-\frac{0.09|X|^2}{4\Delta_*|X|+0.2\Delta_*|X|}\Big)\le\exp\Big(-\frac{|X|}{47\Delta_*}\Big).
\]
Since $\Ex Z-0.3|X|\ge0.33|X|$, we get
$\Prob(Z<0.33|X|)\le\exp(-|X|/(47\Delta_*))$. Outside this event,
$|A_i(W)|\ge Z\ge0.33|X|\ge0.165\,\varepsilon'|W|/L^2\ge g_*|W|$.

For the exponent, $|X|\ge\varepsilon'|W|/(2L^2)$ and $|W|\ge\theta_*u$
give
\[
\frac{|X|}{47\Delta_*}\ge\frac{\varepsilon'\theta_*u}{94\Delta_*L^2}=\frac{2^{19}\ell_*^2Lu}{94\,\Delta_*\rho^2}.
\]
\begin{itemize}
\item If $p_*\le0.11$, then \ref{s3:eqS3} and~\ref{s3:eqS2} give
  $\Delta_*\le3p_*^{-2}\le300\,\ell_*^2\rho^{-2}$, and the exponent is at
  least $2^{19}uL/28200\ge18uL$ (indeed $2^{19}/28200>18.59$; this is the
  tightest inequality of the lemma).
\item If $p_*>0.11$, then \ref{s3:eqS3} gives
  $\Delta_*\le13p_*^{-2}<13/0.0121<1075$, and the exponent is at least
  $2^{19}\ell_*^2uL/(94\cdot1075)\ge5\ell_*^2uL\ge18uL$.
\end{itemize}
The large-$p_*$ regime occurs only for $\rho$ extremely close to $1$. It
is handled directly by the bound $\Delta_*\le13p_*^{-2}$, so no coupling
with a smaller density is needed. This proves~\eqref{s3:eqL17claim}.

\smallskip\noindent\emph{Step 3: all rounds.} For $1\le i\le\ell_*-1$ let
$E_i$ be the event that $|\mathcal B_i|\ge2n/3$ or
$|\mathcal B_{i+1}\setminus\mathcal B_i|\ge g_*|\mathcal B_i|$. Condition
on $V_1,\dots,V_{i-1}$. This fixes $W:=\mathcal B_i\supseteq\mathcal B_1$,
while $V_i$ is still $p_*$-random (by the independence of $V_i$ from
$(V_1,\dots,V_{i-1})$). If $|W|\ge2n/3$ then $E_i$ holds.
Otherwise $A_i(W)\subseteq\mathcal B_{i+1}\setminus\mathcal B_i$ by (O1),
since $A_i(W)\cap W=\emptyset$, and~\eqref{s3:eqL17claim} gives
$\Prob(E_i^c\mid V_1,\dots,V_{i-1})\le e^{-18uL}$. Hence, using $u\ge1$,
\ref{s3:eqS1}, and the fact that $L^3e^{-11L}\le e^{-11}$ for $L\ge1$,
\[
\Prob\Big(\bigcup_{i<\ell_*}E_i^c\Big)\le(\ell_*-1)e^{-18uL}\le2^{10}L^3e^{-11L}\,e^{-7uL}\le0.02\,e^{-7uL}.
\]

\smallskip\noindent\emph{Step 4: growth to two thirds.} Suppose all
$E_i$ hold, and suppose $|\mathcal B_{\ell_*}|<2n/3$. Then
$|\mathcal B_i|<2n/3$ for all $i\le\ell_*$, so
$|\mathcal B_{i+1}|\ge(1+g_*)|\mathcal B_i|$ for $1\le i\le\ell_*-1$, and
$|\mathcal B_{\ell_*}|\ge(1+g_*)^{\ell_*-1}$. Since $g_*\le1/8$, we have
$\ln(1+g_*)\ge g_*-g_*^2/2\ge15g_*/16$. By~\ref{s3:eqS1} and
$\varepsilon'\ge2^{-7}$,
$(\ell_*-1)g_*>(2^{10}L^3-2)\,2^{-10}L^{-2}\ge L-2^{-9}$. Hence
$\ln\big((1+g_*)^{\ell_*-1}\big)\ge\tfrac{15}{16}(L-2^{-9})\ge L\ln2=\ln n$,
so $|\mathcal B_{\ell_*}|\ge n$, a contradiction. Therefore
$|\mathcal B_{\ell_*}|\ge2n/3$.

\smallskip\noindent\emph{Step 5: the last batch.} Condition on
$V_1,\dots,V_{\ell_*-1}$ with $|\mathcal B_{\ell_*}|\ge2n/3$. The set
$V_{\ell_*}$ is still $0.9\rho$-random, so
$|\mathcal B_{\ell_*}\cap V_{\ell_*}|$ is binomial with mean at least
$0.6\rho n$. The lower-tail Chernoff bound for a mean bound $\mu:=0.6\rho n\le\Ex|\mathcal B_{\ell_*}\cap V_{\ell_*}|$
(Cited result~\ref{s1:citChernoffGen}(a)) with $\delta=0.09$ gives
$\Prob(|\mathcal B_{\ell_*}\cap V_{\ell_*}|<0.546\rho n)\le
\exp(-0.0081\cdot0.6\rho n/2)\le\exp(-0.00243\rho n)$. Since $V$ is
$\rho$-random, $|V|\sim\Bin(n,\rho)$, and
$\Prob(|V|>1.09\rho n)\le\exp(-0.0081\rho n/3)=\exp(-0.0027\rho n)$.
Outside both events,
$|\mathcal B_{\ell_*}\cap V_{\ell_*}|\ge0.546\rho n>0.545\rho n\ge|V|/2$,
and then $|\Ball{\ell_*}{G-F}{U,V}|>|V|/2$ by (O2). The probability of
the two events is at most
$2e^{-\rho n/412}\le 2e^{-\rho n/413}e^{-7uL}\le0.01\,e^{-7uL}$. Here we
used Step~0: $7uL\le2^{-36}\rho n\le\rho n(1/412-1/413)$ and
$\rho n>2^{39}$.

\smallskip\noindent\emph{Conclusion.} If
$|\Ball{\ell_*}{G-F}{U,V}|\le|V|/2$, then either some $E_i$ fails, or
(by Step~4) $|\mathcal B_{\ell_*}|\ge2n/3$ and one of the two events of
Step~5 occurs. The total probability is at most
$0.03\,e^{-7uL}\le e^{-7uL}$. The final sentence of the lemma holds
because the event is a function of the random set $V$ alone, and $V$ is
$\rho$-random.
\end{proof}

\begin{lemma}[Proposition 18$^*$ and Lemma 19$^*$]\label{s3:lemP18s}
Let $G$ be an $n$-vertex $(\varepsilon',s_1)$-expander with $n\ge2$ and
$2^{-7}\le\varepsilon'\le1$, let $\rho\in(0,1]$, take the parameters
of~\eqref{s3:eqStar} with $L:=\log n$, and assume
$s_1\ge\theta_*+1$.
\begin{itemize}
\item[(i)] (Proposition 18$^*$.) Every $U\subseteq V(G)$ with
  $1\le|U|\le2n/3$ contains a set $U'$ with
  $|\Nbr_G(U')|\ge\theta_*|U'|$ and
  $|U'|\ge\varepsilon'|U|/(3\theta_*L^2)$.
\item[(ii)] (Lemma 19$^*$.) Let $V$ be a $\rho$-random subset of
  $V(G)$. With probability at least $1-n^{-6}$ the following holds: for
  every nonempty $U\subseteq V(G)$ and every $F\subseteq E(G)$ with
  $|F|\le\mu_*|U|$,
  \[
  |\Ball{\ell_*}{G-F}{U,V}|>|V|/2 .
  \]
\end{itemize}
The hypothesis is $s_1\ge\theta_*+1$ and not $s_1\ge\theta_*$. The
maximality step of \cite[Proposition~18]{BM} (Cited result~\ref{s1:citProp18})
produces vertices with \emph{fewer than $\theta_*+1$} neighbours outside,
so the edge set $F$ built there satisfies $|F|<(\theta_*+1)|U|$.
\deps{\ref{s1:citProp18}, \ref{s1:citLem19}, \ref{s3:lemL17s}, \ref{s1:citDef11}}
\fixes{\RTb{I12}}
\end{lemma}

\begin{proof}
By~\ref{s3:eqS4}, $\theta_*\ge 8d_*\lambda_*$, so $s_1\ge 8d_*\lambda_*$
and Lemma~\ref{s3:lemL17s} applies to $G$. Also $\theta_*\ge1$.

(i) We follow the proof of \cite[Proposition~18]{BM}. Let $U'\subseteq U$
be maximal subject to $|\Nbr_G(U')|\ge\theta_*|U'|$; the empty set
qualifies. If $U'=U$ we are done, since
$\varepsilon'/(3\theta_*L^2)\le1$. Otherwise let $v\in U\setminus U'$, and
let $a_v$ be the number of neighbours of $v$ outside
$U'\cup\Nbr_G(U')$. Then
\[
|\Nbr_G(U'\cup\{v\})|=|\Nbr_G(U')|-\ind{v\in\Nbr_G(U')}+a_v .
\]
By maximality this is less than $\theta_*(|U'|+1)$. Two consequences
follow:
\begin{itemize}
\item $a_v<\theta_*(|U'|+1)-|\Nbr_G(U')|+1\le\theta_*+1$;
\item $|\Nbr_G(U')|<\theta_*(|U'|+1)+1$.
\end{itemize}
Let $F$ be the set of edges $vx$ with $v\in U\setminus U'$ and
$x\notin U'\cup\Nbr_G(U')$. Then
$|F|\le\sum_{v\in U\setminus U'}a_v<(\theta_*+1)|U|\le s_1|U|$.
Every vertex $x\notin U$ with a neighbour in $U$ in $G-F$ lies in
$\Nbr_G(U')$: if the neighbour is in $U'$ this is clear, and if it is in
$U\setminus U'$ then, since the edge is not in $F$,
$x\in U'\cup\Nbr_G(U')$ with $x\notin U'$. By the expansion of $G$
(Cited result~\ref{s1:citDef11}),
\[
\varepsilon'|U|/L^2\le|\Nbr_{G-F}(U)|\le|\Nbr_G(U')|<\theta_*(|U'|+1)+1 .
\]
If $U'=\emptyset$, then the middle term is $0$, which is impossible
because $|U|\ge1$. So $|U'|\ge1$, and then
$\theta_*(|U'|+1)+1\le3\theta_*|U'|$ since $\theta_*\ge1$. Hence
$|U'|\ge\varepsilon'|U|/(3\theta_*L^2)$.

(ii) We follow the proof of \cite[Lemma~19]{BM}
(Cited result~\ref{s1:citLem19}).

\emph{The union bound.} Consider all pairs $(U',F')$ with $U'$
well-expanding and nonempty, $F'\subseteq E(G)$ and $|F'|\le|U'|$. Put
$e_G:=|E(G)|\le n^2/2$; since $G$ has minimum degree $>s_1\ge1$,
$e_G\ge2$. The number of pairs with $|U'|=u$ is at most
$n^u\sum_{f=0}^{u}e_G^f\le n^u\cdot2e_G^u\le n^{3u}$. By
Lemma~\ref{s3:lemL17s}, each pair fails
$|\Ball{\ell_*}{G-F'}{U',V}|>|V|/2$ with probability at most
$e^{-7uL}=n^{-7u\log e}\le n^{-10.09u}$. Hence, with probability at least
\[
1-\sum_{u\ge1}n^{3u}n^{-10.09u}\ \ge\ 1-\frac{n^{-7.09}}{1-n^{-7.09}}\ \ge\ 1-n^{-6},
\]
every such pair satisfies the ball bound. Fix an outcome in this event.

\emph{Deterministic consequence.} Let $U$ be nonempty, and let
$F\subseteq E(G)$.
\begin{itemize}
\item \emph{Case $|U|\le2n/3$ and $|F|\le2\mu_*|U|$.} By (i) there is a
  well-expanding $U'\subseteq U$ with
  $|U'|\ge\varepsilon'|U|/(3\theta_*L^2)=2\mu_*|U|\ge|F|$. The ball
  $\Ball{\ell_*}{G-F}{\cdot,V}$ is monotone in its first argument, so
  $|\Ball{\ell_*}{G-F}{U,V}|\ge|\Ball{\ell_*}{G-F}{U',V}|>|V|/2$.
\item \emph{Case $|U|>2n/3$ and $|F|\le\mu_*|U|$.} Choose
  $\bar U\subseteq U$ with $n/2\le|\bar U|\le2n/3$. This is possible
  for every $n\ge2$: the interval $[n/2,2n/3]$ contains an integer, as
  one checks directly for $n\le5$, and for $n\ge6$ it has length at
  least $1$. Then $|F|\le\mu_*n\le2\mu_*|\bar U|$, and the previous case
  gives $|\Ball{\ell_*}{G-F}{U,V}|\ge|\Ball{\ell_*}{G-F}{\bar U,V}|>|V|/2$.
\end{itemize}
In both cases the conclusion of (ii) holds.
\end{proof}

\begin{lemma}[Lemma 9$_\rho$: multiset linking from ball expansion]\label{s3:lemL9rho}
Let $G$ be an $n$-vertex graph with $n\ge2^{30}$, put $L:=\log n$, and
let $t\ge1$ and $\rho\in(0,1]$. Take $\ell_*$, $\bar s_*$ and
$M_*=\lceil2.1/\rho\rceil$ from~\eqref{s3:eqStar}. Let $V\subseteq V(G)$
with $|V|\ge\rho n/2$, and suppose $\rho n\ge84$. Then
$|V|\ge n/M_*+2$. Suppose that
\[
|\Ball{\ell_*}{G-F}{U,V}|>|V|/2
\]
for every nonempty $U\subseteq V(G)$ and every $F\subseteq E(G)$ with
$|F|\le\bar s_*|U|$. Then $G$ is $(2^{12}L^4,t)$-path connected through
$V$ in the multiset sense of \cite[Definition~7]{BM}
(Cited result~\ref{s1:citDef7}). That is, every multiset of pairs of distinct
vertices of $G$ in which each vertex lies in at most $t$ pairs can be
joined by pairwise edge-disjoint paths of length at most $2^{12}L^4$
whose interior vertices lie in $V$.
\deps{\ref{s1:citLem9}, \ref{s1:citProp8}, \ref{s1:citHaxell}, \ref{s1:citDef7}, \ref{s3:remMultiset}}
\end{lemma}

\begin{proof}
First, $n/M_*\le\rho n/2.1$ and
$\rho n/2-\rho n/2.1=\rho n/42\ge2$, so $|V|\ge n/M_*+2$. Next,
$1\le\ell_*\le2^{10}L^3\le n$ for $n\ge2^{30}$. We re-run the proof of
\cite[Lemma~9]{BM} (Cited result~\ref{s1:citLem9}) with its constants changed,
and with Haxell's theorem (Cited result~\ref{s1:citHaxell}) in place of the
Aharoni--Haxell theorem \cite[Theorem~6]{BM}.
It uses \cite[Proposition~8]{BM} (Cited result~\ref{s1:citProp8}): if
$1\le\ell,t_0\le n$, $|V|\ge4t_0-2$, and every set of $t_0$ vertices has
a ball of radius $\ell$ in $V$ of size greater than $|V|/2$, then among
any $2t_0-1$ pairs of vertices, all $4t_0-2$ of them distinct, some pair
is joined by a path through $V$ of length at most $4\ell\log n$.

Let $\mathcal P=(\{x_i,y_i\})_{i\in[r]}$ be a multiset of pairs of
distinct vertices, indexed by $[r]$, in which every vertex lies in at
most $t$ pairs (counted with multiplicity; see
Remark~\ref{s3:remMultiset}). Put $h:=\lfloor4\ell_*L\rfloor$; then
$h\ge1$, since $\ell_*\ge1$ and $L\ge1$. For $i\in[r]$, let $\mathcal Y_i$
be the set of $x_iy_i$-paths in $G$ that have length at most $h$ and all
interior vertices in $V$. Every $P\in\mathcal Y_i$ has at most $h$ edges.

\emph{Claim: for every nonempty $I\subseteq[r]$ and every $F\subseteq E(G)$
with $|F|<h^2|I|$, there are $j\in I$ and $P\in\mathcal Y_j$ with
$E(P)\cap F=\emptyset$.}
Let $I'\subseteq I$ be maximal such that the pairs $\{x_i,y_i\}$,
$i\in I'$, are pairwise vertex-disjoint; $I'\ne\emptyset$ because
$I\ne\emptyset$. By maximality every pair
indexed by $I$ contains one of the $2|I'|$ vertices covered by $I'$, and
each vertex lies in at most $t$ pairs. Hence $|I|\le2t|I'|$. Also
$|I'|\le n/2$. Put $t_{0*}:=\lceil|I'|/(2M_*)\rceil$. Then:
\begin{itemize}
\item $1\le t_{0*}$ and $2t_{0*}-1\le|I'|$ (if $|I'|\le2M_*$ then
  $t_{0*}=1$; otherwise $2t_{0*}-1\le|I'|/M_*+1\le|I'|$, as $M_*\ge3$);
\item $4t_{0*}-2\le n/M_*+2\le|V|$;
\item $t_{0*}\le n$;
\item $|F|<h^2|I|\le16(\ell_*L)^2\cdot2t|I'|\le64M_*t\,t_{0*}(\ell_*L)^2$.
\end{itemize}
Since $M_*\le2.1/\rho+1\le3.1/\rho$ and $(\ell_*L)^2\le2^{20}L^8$, we get
$64M_*t(\ell_*L)^2\le2^{28}tL^8/\rho=\bar s_*$. Hence $|F|<\bar s_*t_{0*}$.
So every set $U$ of $t_{0*}$ vertices satisfies $|F|\le\bar s_*|U|$, and
therefore $|\Ball{\ell_*}{G-F}{U,V}|>|V|/2$ by hypothesis.

Choose $2t_{0*}-1$ indices of $I'$; their $4t_{0*}-2$ endpoints are
distinct. Apply Cited result~\ref{s1:citProp8} to $G-F$, with $\ell:=\ell_*$
and $t_0:=t_{0*}$. It gives, for one of these indices $j$, an
$x_jy_j$-path $P$ in $G-F$ through $V$ of length at most $4\ell_*L$, hence
at most $h$, as lengths are integers. So $j\in I'\subseteq I$,
$P\in\mathcal Y_j$ and $E(P)\cap F=\emptyset$. This proves the claim.

Now let $\mathcal H$ be the hypergraph with vertex set $[r]\sqcup E(G)$ (a
disjoint union) whose edges are the sets $\{i\}\cup E(P)$ with $i\in[r]$
and $P\in\mathcal Y_i$. Every edge of $\mathcal H$ meets $[r]$ in exactly
one vertex and $E(G)$ in at most $h$ vertices. We apply Cited
result~\ref{s1:citHaxell} to $\mathcal H$, with $[r]$ in the role of its
class $\mathcal X$, $E(G)$ in the role of its other class, and with $h\ge1$ in
place of its parameter $q$. We check its hypothesis. Let $I\subseteq[r]$ and $Z\subseteq E(G)$
with $|Z|\le(2h-1)(|I|-1)$. Then $I\ne\emptyset$, since otherwise
$|Z|\le1-2h<0$. As $h^2-(2h-1)=(h-1)^2\ge0$ and $|I|-1\ge0$,
\[
|Z|\le(2h-1)(|I|-1)\le h^2(|I|-1)<h^2|I| .
\]
So the claim, with $F:=Z$, gives $j\in I$ and $P\in\mathcal Y_j$ with
$E(P)\cap Z=\emptyset$. The edge $e:=\{j\}\cup E(P)$ of $\mathcal H$
satisfies $e\cap[r]=\{j\}\subseteq I$ and, as $j\notin E(G)\supseteq Z$,
$e\cap Z=E(P)\cap Z=\emptyset$. (The claim covers every $Z$ with
$|Z|<h^2|I|$, which is more than the hypothesis of Cited
result~\ref{s1:citHaxell} requires.) By Cited result~\ref{s1:citHaxell},
$\mathcal H$ has a set of pairwise disjoint edges such that every
$i\in[r]$ lies in one of them. For $i\in[r]$ let $e_i$ be an edge of this
set containing $i$. By construction $e_i=\{k\}\cup E(Q)$ for some
$k\in[r]$ and $Q\in\mathcal Y_k$, so $e_i\cap[r]=\{k\}$; as
$i\in e_i\cap[r]$, we get $k=i$, and we put $P_i:=Q$, so that
$E(P_i)\subseteq e_i$. For $i\ne i'$ we have
$e_i\cap[r]=\{i\}\ne\{i'\}=e_{i'}\cap[r]$, so $e_i\ne e_{i'}$, and
therefore $e_i\cap e_{i'}=\emptyset$, as both lie in the set of pairwise
disjoint edges. Hence the paths $P_i$, $i\in[r]$, are pairwise
edge-disjoint $x_iy_i$-paths through $V$ of length at most
$h\le4\ell_*L\le2^{12}L^4$. Repeated pairs receive distinct paths,
because the paths are indexed by $i$.
\end{proof}

\begin{theorem}[Theorem 16$^*$: linking through random sets; for-all multiset form]\label{s3:thmT16s}
Let $X$ be an $N$-vertex $(\varepsilon',s)$-expander with
$\varepsilon'\in[2^{-7},1]$, and put $L:=\log N$. Let $\rho\in(0,1]$ and
$t\ge1$, and let $V$ be a $\rho$-random subset of $V(X)$. Here $X$ is
fixed; if $X$ was itself produced at random, $V$ is independent of that
randomness. Suppose that
\[
\rho N\ge L^2\qquad\text{and}\qquad s\ge 2^{135}\,t\,L^{28}\rho^{-5}.
\]
Then, with probability at least $1-2^{86}\,t\,L^{19}\rho^{-3}N^{-3}$, the
graph $X$ is $(2^{12}L^4,t)$-path connected through $V$. That is, for
\emph{every} multiset $\mathcal P$ of pairs of distinct vertices of $X$
(endpoints anywhere, in $V$ or not) in which each vertex lies in at most
$t$ pairs, there are pairwise edge-disjoint paths of $X$, one for each
pair of $\mathcal P$, each joining its pair, of length at most
$2^{12}L^4$, and with all interior vertices in $V$. By
Lemma~\ref{s3:lemMonotone}(iii), $\mathcal P$ may be chosen after $V$ is
revealed. Moreover:
\begin{itemize}
\item[(a)] The constant $2^{86}$ in the failure probability is absolute.
  No lower bound on $N$ is needed.
\item[(b)] If $N\ge2$, the hypotheses force
  $\rho N>2^{27}L^{5.6}N^{4/5}$, because $s<N$. In particular the
  hypothesis $\rho N\ge L^2$ is implied by the others; it is kept only
  for comparison with \cite[Theorem~16]{BM}. The last batch of the
  sprinkling argument needs $\rho N$ large compared with $L$ times the
  size of the well-expanding sets involved. This is automatic (Step~0 of
  the proof of Lemma~\ref{s3:lemL17s}), so the final-stage union bound
  needs no further hypothesis.
\item[(c)] The requirement on $s$ is exactly linear in $t$ and uniform in
  $\rho\in(0,1]$. Apart from Step~0 of the proof below, which uses the full hypothesis on $s$ to
  obtain (b), $N\ge2^{30}$ and $\rho N\ge84$ (from (b), Step~5 then draws
  $e^{-\rho N/8}\le N^{-3}$),
  the proof uses only $s\ge2K_*(\theta_*+1)$, and
  $2K_*(\theta_*+1)\le2^{130.6}\,tL^{28}\rho^{-5}$ by~\ref{s3:eqS5}. No
  coupling with a smaller density is used, also when $\rho$ is close to
  $1$.
\end{itemize}
\deps{\ref{s3:lemL15p}, \ref{s3:propP13s}, \ref{s3:lemL17s}, \ref{s3:lemP18s}, \ref{s3:lemL9rho}, \ref{s3:lemMonotone}, \ref{s1:citDef11}, \ref{s1:citChernoff}}
\fixes{\RTb{I12}}
\end{theorem}

\begin{proof}
The proof is an assembly, following the proof of \cite[Theorem~16]{BM}
(Cited result~\ref{s1:citThm16}). The loss-free splitting
Lemma~\ref{s3:lemL15p} takes the place of \cite[Lemma~15]{BM}.

\smallskip\noindent\emph{Step 0: $N$ and $\rho N$ are large.}
If $N=1$ there are no pairs of distinct vertices, and there is nothing
to prove. Let $N\ge2$, so $L\ge1$. Taking $U=\{v\}$ in
Cited result~\ref{s1:citDef11} shows that $X$ has minimum degree greater than
$s$. Hence $N>s\ge2^{135}tL^{28}\rho^{-5}\ge2^{135}$. Then
$\rho^5N>\rho^5s\ge2^{135}L^{28}$, so $\rho>2^{27}L^{5.6}N^{-1/5}$, which
is (b). In particular $\rho N\ge84$ and $N\ge2^{30}$. Take the
parameters~\eqref{s3:eqStar} with $n:=N$ and the given $\varepsilon'$,
$\rho$ and $t$.

\smallskip\noindent\emph{Step 1: splitting.} Colour $E(X)$ with $K_*$
colours, uniformly and independently of $V$, and let
$X_1,\dots,X_{K_*}$ be the colour classes, each with vertex set $V(X)$.
By~\ref{s3:eqS5}, $s\ge2K_*(\theta_*+1)\ge40K_*L$. So by
Lemma~\ref{s3:lemL15p}, with probability at least $1-2K_*N^{-5}$, every
$X_i$ is an $(\varepsilon',s/(2K_*))$-expander, and
$s/(2K_*)\ge\theta_*+1$.

\smallskip\noindent\emph{Step 2: balls in each class.} Condition on a
colouring in which all classes are such expanders; $V$ is still
$\rho$-random. Each $X_i$ has $N$ vertices and parameter $\varepsilon'$,
so the parameters~\eqref{s3:eqStar} for $X_i$ are those for $X$. By
Lemma~\ref{s3:lemP18s}(ii) with $s_1:=s/(2K_*)$ and a union bound over
$i$, with probability at least $1-K_*N^{-6}$ the following holds for
every $i$: for every nonempty $U$ and every $F_i\subseteq E(X_i)$ with
$|F_i|\le\mu_*|U|$, we have $|\Ball{\ell_*}{X_i-F_i}{U,V}|>|V|/2$.

\smallskip\noindent\emph{Step 3: the size of $V$.} Since
$|V|\sim\Bin(N,\rho)$, the Chernoff bound (Cited result~\ref{s1:citChernoff})
gives $\Prob(|V|<\rho N/2)\le e^{-\rho N/8}$.

\smallskip\noindent\emph{Step 4: assembly.} Assume the events of
Steps~1--3. Let $U$ be nonempty and let $F\subseteq E(X)$ with
$|F|\le\bar s_*|U|$. The classes partition $E(X)$, so some $i$ has
$|F\cap E(X_i)|\le|F|/K_*\le\bar s_*|U|/K_*\le\mu_*|U|$
by~\ref{s3:eqS5}. Since $X_i-(F\cap E(X_i))\subseteq X-F$, balls can only
grow when passing to $X-F$, and
\[
|\Ball{\ell_*}{X-F}{U,V}|\ge|\Ball{\ell_*}{X_i-(F\cap E(X_i))}{U,V}|>|V|/2 .
\]
So the hypotheses of Lemma~\ref{s3:lemL9rho} hold for $G:=X$:
$N\ge2^{30}$, $|V|\ge\rho N/2$ and $\rho N\ge84$. Hence $X$ is
$(2^{12}L^4,t)$-path connected through $V$, in the multiset sense.

\smallskip\noindent\emph{Step 5: probability.} The failure probability
is at most $2K_*N^{-5}+K_*N^{-6}+e^{-\rho N/8}$. By (b),
$\rho N/8>2^{24}N^{4/5}\ge3\ln N$, so $e^{-\rho N/8}\le N^{-3}$. By
Lemma~\ref{s3:lemMonotone}(iii), the random colouring of Step~1, which is
auxiliary, does not affect the event, which depends only on $(X,V)$. So
the event fails with probability at most
\[
3K_*N^{-5}+N^{-3}\le\big(3\cdot2^{83.6}+1\big)tL^{19}\rho^{-3}N^{-3}\le2^{86}\,tL^{19}\rho^{-3}N^{-3},
\]
using~\ref{s3:eqS5} and $tL^{19}\rho^{-3}\ge1$. This proves the main
statement and~(a).

For (c): apart from Step~0, whose consequences $N\ge2^{30}$ and $\rho N\ge84$
(used in Step~4) and $e^{-\rho N/8}\le N^{-3}$ (used in Step~5, through (b)) rest
on the full hypothesis $s\ge2^{135}tL^{28}\rho^{-5}$, the requirements on $s$
used above are $s\ge40K_*L$ (Step~1) and $s/(2K_*)\ge\theta_*+1$ (Step~2). Both follow from
$s\ge2K_*(\theta_*+1)$, and $2K_*(\theta_*+1)\le2^{130.6}tL^{28}\rho^{-5}$
by~\ref{s3:eqS5}. Only $\bar s_*$, and hence $K_*$, depends on $t$, and
linearly up to rounding. The parameters~\eqref{s3:eqStar} are defined for
every $\rho\in(0,1]$. The regime $p_*>0.11$, reached only for $\rho$ very
close to $1$, is treated directly in Lemma~\ref{s3:lemL17s} through
$\Delta_*\le13p_*^{-2}$ (fact~\ref{s3:eqS3}).
\end{proof}

%% ----------------------------------------------------------------------
\subsection{Candidate sets of bounded multiplicity}

\begin{lemma}[Lemma HB: candidate sets of bounded multiplicity (the name HB is historical); Hall argument]\label{s3:lemHB}
Let $X$ be an $N$-vertex $(\varepsilon',s)$-expander with $N\ge2$ and
$2^{-7}\le\varepsilon'\le1$, and put $L:=\log N$. Let $m$ be an integer
with $1\le m$ and $2m\le s$, and put $b:=\lceil2mL^2/\varepsilon'\rceil$.
Thus $b=\lceil64L^2m\rceil$, $\lceil128L^2m\rceil$ or
$\lceil256L^2m\rceil$ at $\varepsilon'=2^{-5}$, $2^{-6}$ or $2^{-7}$.
Then every vertex $w$ has a set $\Aw(w)\subseteq\Nbr_X(w)$ with
$|\Aw(w)|=m$, such that every vertex of $X$ lies in at most $b$ of the
sets $\Aw(w)$, $w\in V(X)$.
\deps{\ref{s1:citHall}, \ref{s1:citDef11}}
\end{lemma}

\begin{proof}
For $w\in V(X)$, its degree is $d_X(w)=|\Nbr_X(w)|$. As $N\ge2$, we have
$L\ge1$ and $d_X(w)\ge\delta(X)>s\ge2m$ for every $w$ (Cited
result~\ref{s1:citDef11}); in particular $d_X(w)-m>m\ge1$ (not needed
below). Let
$\vec E(X):=\{(w,u)\in V(X)\times V(X):\ wu\in E(X)\}$ be the set of
ordered pairs along edges (every edge gives two), and for $w\in V(X)$ let
$\vec E_X(w):=\{(w,u):\ u\in\Nbr_X(w)\}$. Then $|\vec E_X(w)|=d_X(w)$, and
distinct elements of $\vec E_X(w)$ have distinct second coordinates.

\emph{The auxiliary bipartite graph.} Its left side is $\vec E(X)$. Its
right side is the set of \emph{slots}
\[
\mathcal B:=\{(u,j,1):\ u\in V(X),\ j\in[b]\}\ \cup\
\{(w,i,2):\ w\in V(X),\ i\in[d_X(w)-m]\} ;
\]
$(u,j,1)$ is a \emph{load slot} of $u$, and $(w,i,2)$ is a \emph{spare
slot} of $w$. A pair $a=(w,u)\in\vec E(X)$ is joined to the $b$ load slots
of its second coordinate $u$ and to the $d_X(w)-m$ spare slots of its first
coordinate $w$:
\[
\mathcal T(a):=\{(u,j,1):\ j\in[b]\}\ \cup\ \{(w,i,2):\ i\in[d_X(w)-m]\}.
\]

\emph{Reduction to Hall's condition.} Suppose that
$\chi:\vec E(X)\to\mathcal B$ is injective and $\chi(a)\in\mathcal T(a)$ for
every $a\in\vec E(X)$. Fix $w\in V(X)$. If $a\in\vec E_X(w)$ and $\chi(a)$ is
a spare slot, then $\chi(a)$ is one of the $d_X(w)-m$ spare slots of $w$; as
$\chi$ is injective, at most $d_X(w)-m$ elements of $\vec E_X(w)$ are of
this kind. So at least $d_X(w)-(d_X(w)-m)=m$ elements
$(w,u)\in\vec E_X(w)$ have $\chi(w,u)$ a load slot, and then it is a load
slot of $u$. Let $\Aw(w)$ be the set of second coordinates of $m$ such
elements (the first $m$ in a fixed ordering of $V(X)$). Then
$\Aw(w)\subseteq\Nbr_X(w)$, and $|\Aw(w)|=m$ because distinct elements of
$\vec E_X(w)$ have distinct second coordinates. Now fix $u\in V(X)$. If
$u\in\Aw(w)$, then $\chi(w,u)$ is one of the $b$ load slots of $u$. The
pairs $(w,u)$ with distinct $w$ are distinct and $\chi$ is injective, so
$u\in\Aw(w)$ holds for at most $b$ vertices $w$. This is the statement. By
Hall's theorem (Cited result~\ref{s1:citHall}, applied to the finite index
set $\vec E(X)$ and the finite sets $\mathcal T(a)$), such a $\chi$ exists
provided
\begin{equation}\label{s3:eqHBhall}
|\mathcal E|\ \le\ \Bigl|\bigcup_{a\in\mathcal E}\mathcal T(a)\Bigr|\qquad\text{for every }\mathcal E\subseteq\vec E(X).
\end{equation}

\emph{Proof of~\eqref{s3:eqHBhall}.} Fix $\mathcal E\subseteq\vec E(X)$.
Let $\mathcal S$ be the set of first coordinates and $\mathcal R$ the set of
second coordinates of the elements of $\mathcal E$. For $w\in\mathcal S$ put
$\mathcal E_w:=\mathcal E\cap\vec E_X(w)$, so that
$1\le|\mathcal E_w|\le d_X(w)$ and
$|\mathcal E|=\sum_{w\in\mathcal S}|\mathcal E_w|$. The union
in~\eqref{s3:eqHBhall} consists of the $b|\mathcal R|$ load slots of the
vertices of $\mathcal R$ and the $\sum_{w\in\mathcal S}(d_X(w)-m)$ spare
slots of the vertices of $\mathcal S$, and all these slots are distinct.
So \eqref{s3:eqHBhall} is equivalent to
\begin{equation}\label{s3:eqHBdef}
\sum_{w\in\mathcal S}\bigl(|\mathcal E_w|-d_X(w)+m\bigr)\ \le\ b|\mathcal R| .
\end{equation}
Let $\mathcal S':=\{w\in\mathcal S:\ d_X(w)-|\mathcal E_w|\le m-1\}$. A
summand of~\eqref{s3:eqHBdef} with $w\in\mathcal S\setminus\mathcal S'$ is
at most $0$, and one with $w\in\mathcal S'$ is at most $m$, as
$|\mathcal E_w|\le d_X(w)$. So the left side of~\eqref{s3:eqHBdef} is at
most $m|\mathcal S'|$, and \eqref{s3:eqHBdef} holds if
$\mathcal S'=\emptyset$ (this includes $\mathcal E=\emptyset$). Assume
$\mathcal S'\ne\emptyset$.

Let $w\in\mathcal S'$ and $u\in\Nbr_X(w)\setminus\mathcal R$. Then
$(w,u)\in\vec E_X(w)\setminus\mathcal E_w$, since otherwise
$u\in\mathcal R$; distinct such $u$ give distinct pairs $(w,u)$. As
$|\vec E_X(w)\setminus\mathcal E_w|=d_X(w)-|\mathcal E_w|\le m-1$, we get
\begin{equation}\label{s3:eqHBout}
|\Nbr_X(w)\setminus\mathcal R|\ \le\ m-1\qquad\text{for every }w\in\mathcal S'.
\end{equation}
If $|\mathcal S'|\le2N/3$, put $\mathcal U:=\mathcal S'$. Otherwise let
$\mathcal U\subseteq\mathcal S'$ be any set with
$|\mathcal U|=\lceil|\mathcal S'|/2\rceil$; then
$1\le|\mathcal U|\le\lceil N/2\rceil\le2N/3$, because $N\ge2$ (for $N=2$,
$\lceil N/2\rceil=1\le4/3$; for $N\ge3$, $\lceil N/2\rceil\le(N+1)/2\le2N/3$).
In both cases
\[
1\le|\mathcal U|\le2N/3\qquad\text{and}\qquad|\mathcal U|\ge|\mathcal S'|/2 .
\]
Let $F$ be the set of edges $wu\in E(X)$ with $w\in\mathcal U$ and
$u\notin\mathcal R$. Every edge of $F$ is of the form $wu$ with
$w\in\mathcal U$ and $u\in\Nbr_X(w)\setminus\mathcal R$, so,
by~\eqref{s3:eqHBout},
$|F|\le\sum_{w\in\mathcal U}|\Nbr_X(w)\setminus\mathcal R|\le(m-1)|\mathcal U|\le s|\mathcal U|$,
since $m-1<2m\le s$. Let $v\in\Nbr_{X-F}(\mathcal U)$: then
$v\notin\mathcal U$ and $vw\in E(X)\setminus F$ for some $w\in\mathcal U$.
As $w\in\mathcal U$ and $vw\notin F$, we have $v\in\mathcal R$. Hence
$\Nbr_{X-F}(\mathcal U)\subseteq\mathcal R$, and the expansion of $X$
(Cited result~\ref{s1:citDef11} with $n=N$, applied to $\mathcal U$ and
$F$) gives
\[
|\mathcal R|\ \ge\ |\Nbr_{X-F}(\mathcal U)|\ \ge\ \frac{\varepsilon'|\mathcal U|}{L^2}\ \ge\ \frac{\varepsilon'|\mathcal S'|}{2L^2}.
\]
Since $b\ge2mL^2/\varepsilon'$, it follows that
\[
b|\mathcal R|\ \ge\ \frac{2mL^2}{\varepsilon'}\cdot\frac{\varepsilon'|\mathcal S'|}{2L^2}\ =\ m|\mathcal S'| ,
\]
which is at least the left side of~\eqref{s3:eqHBdef}. This
proves~\eqref{s3:eqHBhall}, and hence the lemma.
\end{proof}

%% ----------------------------------------------------------------------
\subsection{The lending colouring}

Fix a valid $\HBtp$ run on $G$ (Definition~\ref{s2:defHBtp}). Every
ancestor $Y$ (Definition~\ref{s2:defAncestors}) splits the edges of its
expander $H_Y$ into two halves:
\begin{itemize}
\item an \emph{own} half, which only $Y$'s own vortex device uses;
\item a \emph{lend} half, which is coloured into exclusive classes, each
  lent to exactly one later consumer.
\end{itemize}
The random data of this subsection form part of stage~1 of the
randomness. The other stage-1 data, introduced in later sections, are
drawn independently of them.

\begin{definition}[stage-1 lending data: the two-stage colouring COL-JV and the JS labels]\label{s3:defCOL}
For every ancestor $Y$ of round $r$, with $V(Y)$, $H_Y$,
$\varepsilon_Y$, $s_Y$ and $L_Y=\log|V(Y)|$ as in
Definition~\ref{s2:defAncestors}, the random data (i)--(iv) below are
drawn. The data of distinct ancestors are independent.
\begin{itemize}
\item[(i)] \emph{Own/Lend split.} Every edge of $H_Y$ independently lies
  in $\Own_Y$ or in $\Lend_Y$, with probability $1/2$ each.
\item[(ii)] \emph{Lent classes.} Consider the three index families
  \[
  \begin{aligned}
  \IU(Y)&:=\{(l,c,\uslot):\ r+2\le l\le R,\ c\in[4],\ \uslot\in[\Tslot_Y]\}\ \text{ if $Y$ is light, where } \Tslot_Y:=\lceil L_Y^2\rceil,\\
  \IU(Y)&:=\emptyset\ \text{ if $Y$ is standalone,}\\
  \IJS(Y)&:=\{(l,j):\ r+2\le l\le R,\ 0\le j<\KJS_l\},\ \text{ where } \KJS_l:=M_l^2,\\
  \IJV(Y)&:=\{l:\ r+2\le l\le R\}.
  \end{aligned}
  \]
  These families are regarded as pairwise disjoint (their indices carry
  the tags U, JS, JV). Put
  $\klend(Y):=|\IU(Y)|+|\IJS(Y)|+|\IJV(Y)|$.
  \begin{itemize}
  \item If $r\le R-2$, every edge of $\Lend_Y$ independently receives an
    index chosen uniformly from $\IU(Y)\sqcup\IJS(Y)\sqcup\IJV(Y)$. The
    edges with index $(l,c,\uslot)\in\IU(Y)$ form the \emph{U-lent class}
    $\LU_{Y,l,c,\uslot}$. The edges with index $(l,j)\in\IJS(Y)$ form the
    \emph{JS-lent class} $\LJS_{Y,l,j}$. The edges with index
    $l\in\IJV(Y)$ form the \emph{JV-lent class} $\LJV_{Y,l}$. Every
    class is regarded as a graph with vertex set $V(Y)$.
  \item If $r\ge R-1$, all three families are empty, $\klend(Y)=0$, and
    $\Lend_Y$ has no classes.
  \end{itemize}
\item[(iii)] \emph{Own classes (light $Y$ only).} Put
  $J_Y:=\lfloor\log(L_Y/8)\rfloor$ and $\kown:=4J_Y+1$. Every edge of
  $\Own_Y$ independently receives one of $\kown$ labels, chosen
  uniformly. This yields the own classes $\vR_{j,c}$ ($0\le j<J_Y$,
  $c\in[4]$) and $\vM$ of $Y$, each regarded as a graph on $V(Y)$. They
  are the own classes of the four-phase P-vortex run on $Y$ in later
  sections. For standalone $Y$, $\Own_Y$ is not subdivided at stage~1.
\item[(iv)] \emph{JS labels.} For every $r+2\le l\le R$ and every
  $y\in V(Y)$, independently, draw a label
  $\lab_{Y,l}(y)\in\{\ast\}\cup\{0,1,\dots,\KJS_l-1\}$ with
  $\Prob(\lab_{Y,l}(y)=j)=\rho_l:=M_l^{-4}$ for each $j$, so that
  $\Prob(\lab_{Y,l}(y)=\ast)=1-M_l^{-2}$. Put
  $T_j(Y,l):=\{y\in V(Y):\lab_{Y,l}(y)=j\}$.
\end{itemize}
Formally, every edge $e$ of $H_Y$ carries three independent uniform
random variables (a fair bit, a lent index and an own label), and the
labels of (iv) are independent of all edge variables. Consequently:
\begin{itemize}
\item colours of distinct edges are independent;
\item the colourings of distinct ancestors are independent;
\item the labels are independent of the colours;
\item given $\Lend_Y$, stage (ii) is a uniform $\klend(Y)$-colouring of
  $\Lend_Y$, and given $\Own_Y$, stage (iii) is a uniform
  $\kown$-colouring of $\Own_Y$;
\item for fixed $(l,j)$, the set $T_j(Y,l)$ is a $\rho_l$-random subset
  of $V(Y)$, independent of the colouring, and the sets $T_j(Y,l)$,
  $0\le j<\KJS_l$, are pairwise disjoint.
\end{itemize}

\emph{Derived quantities.} For $\klend(Y)\ge1$ put
$p_Y:=1/(2\klend(Y))$. For each fixed $l$ with $r+2\le l\le R$, this is
the probability that a fixed edge of $H_Y$ lies in $\LJV_{Y,l}$. Also put
$\tJS_l:=2M_l+2$. For a light ancestor $Y$ of round $r$ put
$t_Y:=\lceil\lambda_r^{1.6}\rceil$, the \emph{U-multiplicity parameter}
of $Y$: it is the multiplicity bound $t$ at which the U-lent classes of
$Y$ are path connected (Lemma~\ref{s3:lemCOL}(c)). It depends only on
the round of $Y$, and $t_Y\ge M_l$ for every $r+2\le l\le R$, because
$M_l\le\lambda_{l-2}^{1.6}\le\lambda_r^{1.6}$
(Lemma~\ref{s2:lemTower}(a),(b)).

\emph{Consumers} (property COL(g)). If $r\le R-2$, each part of the data
has exactly one consumer:
\begin{itemize}
\item $\Own_Y$, together with its own classes if $Y$ is light, is used
  only by the vortex device of $Y$ itself. This is the transparent
  P-vortex if $Y$ is standalone. If $Y$ is light, it is the four-phase
  P-vortex, or VX$^+$ if $Y$ is demoted.
\item $\LU_{Y,l,c,\uslot}$ is used only by the Theorem-U chaining of the
  U-bundle $\Bdl_{l,c}$, in slot $\uslot$.
\item $\LJS_{Y,l,j}$ is used only by the JS-LC systems of the pair
  $(Y,l)$, at junction $j$.
\item $\LJV_{Y,l}$ is used only by the round-$l$ J-consumer.
\end{itemize}
Lent edges that their consumer does not use are junk for $Y$'s own
device: they are added to the edge set that this device decomposes.
If $r\ge R-1$, then $\Own_Y$ has the same consumer, and $\Lend_Y$ has no
classes and no other consumer: all of $\Lend_Y$ is junk for $Y$'s own
device. Indeed $\Lent(Y)=\emptyset$, since it collects edges used at
rounds $l\ge r+2>R$, so $\Erem(Y)=E_r(Y)$ for standalone $Y$ and
$H_0(Y)=E_r(Y)$ for light non-demoted $Y$ (Construction~\ref{s6:consOrder});
a demoted $Y$ runs on $E_r(Y)$ itself.
\deps{\ref{s2:defAncestors}, \ref{s2:propStructure}, \ref{s2:lemTower}}
\fixes{\textsf{CR1-PV} (the parameter $t_Y$)}
\end{definition}

The graphs $H_Y$ of all ancestors $Y$ of all rounds are pairwise
edge-disjoint (Proposition~\ref{s2:propStructure}(iii)). So the
colourings of distinct ancestors act on disjoint edge sets. The label
distribution in (iv) is well defined because
$\KJS_l\rho_l=M_l^{-2}\le1$.

\begin{table}[htbp]
\begin{center}
\small
\begin{tabular}{|p{0.55cm}|p{4.6cm}|p{5.7cm}|p{2.4cm}|}
\hline
No. & Requirement (use) & Sufficient inequality in $\lambda=2^\mu$ & $(a_i,b_i,c_i)$\\\hline
1 & $\thGC_r(Z^0)>\tau_r$ (rule GC) & $\lambda^{102.5}>257(\lambda+6\mu+20)^{102}$; crude form $\lambda^{1/2}>2^{111}$ & $(0.5,112,0)$\\\hline
2 & $P_{l-2}\ge M_l^{13}$, with $\lambda:=\lambda_{l-2}$ & $\lambda^{103}\ge\bar M^{13}$ & $(103,0,26A)$\\\hline
3 & Lemma 15$^+$, lent level: $s_Y/4\ge40\,\klend L_Y$ & $\lambda^{99}\ge640\,\bar k$ & $(96,19,4A)$\\\hline
4 & T16$^*$ for U-lent classes: $s_Y/(8\klend)\ge2^{135}\,t_Y\,L_Y^{28}(12L_Y^5)^5$, $t_Y=\lceil\lambda^{1.6}\rceil$ & $\lambda^{42.1}\ge2^{193}\cdot12^5$ (with row~9)$^{\dagger}$ & $(42.1,211,0)$\\\hline
5 & T16$^*$ for JS-lent classes: $s_Y/(8\klend)\ge2^{135}(3M)L_Y^{28}M^{20}$ & $\lambda^{72}\ge3\cdot2^{167}\cdot\bar k\,\bar M^{21}$ & $(69,178,46A)$\\\hline
6 & T16$^*$ failures over $\IJS(Y)$ sum to at most $|V(Y)|^{-2}/4$ & $\lambda^{84}\ge2^{110}\bar M^{15}$ & $(84,110,30A)$\\\hline
7 & T16$^*$ failures over $\IU(Y)$ sum to at most $|V(Y)|^{-2}/4$ & $\lambda^{64.4}\ge2^{127}\cdot12^4$ & $(64.4,142,0)$\\\hline
8 & own devices: (a) $s_r/4\ge2^{150}L_Y^{42}$; (b) $s_r/4\ge2^{146}L_Y^{38}\log L_Y$ (the threshold of VX$^+$ for standalone parts; not used, as VX$^+$ is applied only to demoted light parts, Lemma~\ref{s5:lemDemoted}; cf.\ Remark~\ref{s4:remTPVVX}); (c) $s_r/2\ge2^{151}L_Y^{38}\log L_Y$; (d) $s':=s_r/(16\kown)\ge2^{145}L_Y^{41}$ and $s'\ge2\lceil L_Y^6\rceil$ & (a) $\lambda^{58}\ge2^{194}$; (b) $\lambda^{62}\ge2^{186}(\mu+1)$; (c) $\lambda^{62}\ge2^{190}(\mu+1)$; (d) $\lambda^{58}\ge2^{191}$ and $\lambda^{99}\ge64((2\lambda)^6+1)$ & $(58,194,1)$\\\hline
9 & $\klend(Y)\le\lambda^{3.3}$ & $\bar k\le\lambda^{3.3}$ & $(0.3,9,4A)$\\\hline
10 & $p_Y\ge\lambda^{-4}$ & $2\bar k\le\lambda^{4}$ & $(1,10,4A)$\\\hline
11 & $\lambda_{l-2}^{95}/(8M_l^2)\ge2^{10}M_l^{10}$, with $\lambda:=\lambda_{l-2}$ & $\lambda^{95}\ge2^{13}\bar M^{12}$ & $(95,13,24A)$\\\hline
12 & $M_l\le\lambda_{l-2}^{1.6}$, with $\lambda:=\lambda_{l-2}$ & $\bar M\le\lambda^{1.6}$ & $(1.6,0,2A)$\\\hline
13 & \Gstar & $\lambda^{36}\ge2^{240}(A\mu)^{46A}$ & $(36,240,46A)$\\\hline
\end{tabular}
\end{center}
\caption{The COL-JV table (Lemma~COL-JV below). Here $\lambda:=\lambda_r$
and $\mu:=\log\lambda$, except in rows 2, 11 and 12, where
$\lambda:=\lambda_{l-2}$ for a round $3\le l\le R$. We write
$\bar M:=(A\mu)^{2A}$, with $A=105$, and $\bar k:=192\lambda^3+\frac43\bar M^2$.
Column~3 is a sufficient form of the requirement, obtained from the
standing bounds (B1)--(B7) of the proof. Column~4 gives, for row~$i$, reals
$a_i>0$ and $b_i,c_i\ge0$ such that, for $\mu\ge6$, the inequality
$a_i\mu-b_i-c_i\log(A\mu)\ge0$ implies every inequality in column~3 of row~$i$
(Lemma~\ref{s3:lemCOLJVev}, which has no hypothesis); hence every
column-3 inequality holds for all sufficiently large $\mu$, and \Gc{1}(f)
requires it for every $\mu\ge\log\log\Dstar$. The inequalities
$a_i\mu-b_i-c_i\log(A\mu)\ge0$ defined by column~4 are cruder than column~3; they
are not part of \Gc{1}, are not asserted in Lemma~COL-JV, and may fail at some
$\mu\ge\log\log\Dstar$ although \Gc{1} holds. No galactic threshold is
evaluated: the thresholds of Lemma~\ref{s3:lemCOLJVev} are explicit, but their
numerical values are never computed. Rows for off-route devices are omitted.
$^{\dagger}$Row~4: column~3 implies the column-2 requirement together with
row~9.}
\label{s3:tabCOLJV}
\end{table}

\begin{lemma}[the COL-JV table; \RTb{I10}]\label{s3:lemCOLJV}
Assume condition~\ref{s1:condG1}. Let $Y$ be an ancestor of round $r\le R$, and write
$\lambda:=\lambda_r$ and $\mu:=\log\lambda$; if $r\le R-2$, write also $M:=M_{r+2}$. Then:
\begin{itemize}
\item[(i)] If $r\ge R-1$, then $\klend(Y)=0$. If $r\le R-2$, then
  $|\IU(Y)|\le12L_Y^3$, $|\IJS(Y)|\le\frac43M^2$ and
  $|\IJV(Y)|=R-r-1\le2\logs d_r+1$. Hence
  \[
  \klend(Y)\le12L_Y^3+\tfrac43M^2+2\logs d_r+2\le24L_Y^3+\tfrac43M^2\le\bar k\le\lambda^{3.3},
  \]
  and consequently $p_Y\ge\lambda^{-4}$. The JV family has
  $R-r-1\le2\logs d_r+1$ members. Its size is not bounded by an
  absolute constant, since $d_1$ may be of order $n$; the count
  ``$+16$'' must not be used.
\item[(ii)] Every requirement in column~2 and every inequality in column~3
  of Table~\ref{s3:tabCOLJV} holds. (Column~4 is not part of this assertion:
  its triples serve only Lemma~\ref{s3:lemCOLJVev}, and the inequalities
  $a_i\mu-b_i-c_i\log(A\mu)\ge0$ that they define need not hold at $\mu$.) More
  precisely, for each row every inequality in column~3 holds at every
  $\mu\ge\log\log\Dstar$ by \Gc{1}(f); in particular it holds at
  $\mu=\log\lambda_r$ and, for rows~2, 11 and~12, also at
  $\mu=\log\lambda_{l-2}$ for every round $3\le l\le R$, both of which are
  at least $\log\log\Dstar$. At the value of $\lambda$ fixed in the caption of
  Table~\ref{s3:tabCOLJV}, these inequalities imply the requirement in
  column~2 (for row~4, together with row~9). So every requirement holds, in
  every round.
  The scope of the rows is as follows.
  \begin{itemize}
  \item Rows~1, 3 and~8 (rule GC, the three levels of Lemma~15$^+$, and
    the own-device thresholds) hold for every round $r\le R$, that is,
    also for $r\in\{R-1,R\}$: row~1 for every round-$r$ pre-part, and
    rows~3 and~8 for every ancestor $Y$ of round $r$. They involve
    neither $M_{r+2}$ nor any lent class; for $r\ge R-1$ the lent level
    of row~3 is void, because $Y$ has no lent classes.
  \item Rows~4--7 and~10 concern lent classes or $p_Y$. They are
    asserted for $r\le R-2$; for $r\ge R-1$ there are no lent classes
    and $p_Y$ is not defined. Row~9 holds for every $r\le R$, trivially
    when $\klend(Y)=0$.
  \item Rows~2, 11 and~12 concern $\lambda_{l-2}$ for rounds
    $3\le l\le R$, and row~13 is \Gstar; they do not depend on $Y$.
  \end{itemize}
\end{itemize}
\deps{\ref{s1:condG1}, \ref{s1:condGstar}, \ref{s2:lemTower}, \ref{s3:defCOL}, \ref{s2:defAncestors}, \ref{s2:propStructure}, \ref{s3:lemL15p}, \ref{s3:thmT16s}}
\fixes{\RTb{I10}, \RTb{I11}, \textsf{CR1-PV} (rows~4 and~7 with $t_Y$)}
\end{lemma}

\begin{proof}
Every round $r\le R$ has $d_r\ge\Dstar$, so $\mu=\log\lambda_r\ge\log\log\Dstar$, and
the items of \Gc{1} hold at $\mu$. The same holds for $\lambda_{l-2}$ with
$3\le l\le R$. We use
the following standing bounds.
\begin{itemize}
\item[(B1)] $L_Y\le\log M_r\le2\lambda$ (Lemma~\ref{s2:lemTower}(b)).
\item[(B2)] If $r\le R-2$: $M_l\le M\le\bar M$ for $r+2\le l\le R$, and
  $M_{l+1}\le M_l/2$ for $r+2\le l<R$ (Lemma~\ref{s2:lemTower}(b)). Likewise
  $M_l\le(A\log\lambda_{l-2})^{2A}$ for every $3\le l\le R$.
\item[(B3)] $s_r\ge\lambda^{100}$ and $s_r\le2\Lambda_r^{100}$ (Lemma~\ref{s2:lemTower}(b)),
  and $s_Y\ge s_r/2$ (Definition~\ref{s2:defAncestors}).
\item[(B4)] $|V(Y)|\ge P_r/2\ge\lambda^{103}/2$. A standalone $Y$ has
  $|V(Y)|=|Y^0|\ge P_r$, and a light $Y$ has $|Y|\ge|Y^0|/2\ge P_r/2$
  (Proposition~\ref{s2:propStructure}(iv)); and $P_r=\lceil\lambda^{C'}\rceil$
  with $C'=103$.
\item[(B5)] $R-r\le2\logs d_r+2$ and $2^{2\logs d_r+2}\le\lambda$
  (Lemma~\ref{s2:lemTower}(a),(d)), so $2\logs d_r+2\le\mu$.
\item[(B6)] $\Lambda_r\le\lambda+6\mu+20$. Indeed, with
  $T:=\tHB_rd_r=\lambda^2\,2^{\lambda}$ we have $\log T=\lambda+2\mu$. Put
  $B:=2^{16}T\log^4T$. As $\lambda=2^\mu\ge2^{256}$ ($\mu\ge2^8$ by \Gc{1}(a)) and
  $\log T=\lambda+2\mu\ge1$, we have $B\ge2^{16}T\ge2^{16+\lambda}>2^{161}\ge2^{40}$,
  so $M_r\le B+1$ by (R2) (Definition~\ref{s2:defHBtp}), and
  $\Lambda_r=\log M_r\le\log B+\log(1+1/B)\le\log B+1/(B\ln2)\le\log B+2^{-160}$,
  since $1/(B\ln2)<2^{-161}/0.69<2^{-160}$. Next
  $\log B=16+\lambda+2\mu+4\log(\lambda+2\mu)$, and
  $\log(\lambda+2\mu)=\mu+\log(1+2\mu/\lambda)\le\mu+2\mu/(\lambda\ln2)\le\mu+0.0025$,
  because $\mu\ge2^8$ and $2\mu/\lambda=2\mu2^{-\mu}\le2^{-240}$. Hence
  $\Lambda_r\le\lambda+6\mu+16.01+2^{-160}\le\lambda+6\mu+20$.
\item[(B7)] If $r\le R-2$ and $r+2\le l\le R$: $\tJS_l=2M_l+2\le3M_l\le3M$
  (as $M_l\ge2^{40}$), and $\rho_l^{-1}=M_l^4\le M^4$.
\end{itemize}
Only (B2) and (B7) involve $M$, and they are used only for $r\le R-2$.
(B1) and (B3)--(B6) hold for every round $r\le R$.

\emph{(i) Counting.} If $r\ge R-1$, the three index families are empty
(Definition~\ref{s3:defCOL}(ii)), so $\klend(Y)=0$. Let now $r\le R-2$.
\begin{itemize}
\item $|\IU(Y)|=4\Tslot_Y(R-r-1)\le4(L_Y^2+1)(2\logs d_r+1)$. By (B5)
  and (B4), $2\logs d_r+1<\mu\le L_Y$, since $L_Y\ge103\mu-1$. Hence
  $|\IU(Y)|\le4(L_Y^2+1)L_Y\le12L_Y^3$.
\item $|\IJS(Y)|=\sum_{l=r+2}^RM_l^2\le M^2\sum_{i\ge0}4^{-i}=\frac43M^2$
  by (B2).
\item $|\IJV(Y)|=R-r-1\le2\logs d_r+1$ by (B5).
\end{itemize}
The second inequality of the display holds because $2\logs d_r+2\le\mu\le12L_Y^3$.
The third follows from (B1) and (B2). The fourth is the sufficient
inequality of row~9, verified in (ii) below.
Then $p_Y=1/(2\klend(Y))\ge1/(2\lambda^{3.3})\ge\lambda^{-4}$, and row~10
gives the same conclusion directly.

\emph{(ii) The rows.} In each row, the passage from column~2 to column~3
uses only (B1)--(B7) and the fact that the requirement is monotone in
$L_Y$, $M_l$, $\klend$, $s_r$ and $|V(Y)|$ in the appropriate direction.
Rows~1, 3 and~8 use only (B1), (B3)--(B6), Lemma~\ref{s2:lemTower}(c)
and the bound $\klend(Y)\le\bar k$ of (i), which is trivial when
$\klend(Y)=0$; so they hold for every round $r\le R$. Rows~4--7, 9 and~10
are derived for $r\le R-2$, where (B2) and (B7) are available; for
$r\ge R-1$, row~9 reads $0\le\lambda^{3.3}$.
We now give the derivations.
\begin{itemize}
\item[1.] For a round-$r$ pre-part $Z^0$ we have $|Z^0|\ge P_r\ge\lambda^{103}$, so
  $\thGC_r(Z^0)=\lceil|Z^0|\lambda^{-1/2}\rceil\ge\lambda^{102.5}$,
  and $\tau_r\le128s_r\Lambda_r^2+1\le257\Lambda_r^{102}$ by (B3). Now
  use (B6). The crude form uses $\tau_r\le2^{111}\lambda^{102}$
  (Lemma~\ref{s2:lemTower}(c)).
\item[2.] $P_{l-2}\ge\lambda_{l-2}^{103}$, and $M_l\le\bar M$ by (B2).
\item[3.] $s_Y/4\ge\lambda^{100}/8$ and $40\klend L_Y\le80\lambda\bar k$.
  The other two levels of Lemma~\ref{s3:lemL15p} are $s_Y\ge80L_Y$ (the
  Own/Lend split, $k=2$) and $s_r/8\ge40\kown L_Y$ (own classes). They
  follow from $\lambda^{100}/2\ge160\lambda$ and
  $\lambda^{100}/8\ge160\lambda^2$, using $\kown\le L_Y\le2\lambda$, and
  both hold for every $\mu\ge1$.
\item[4.] U-lent classes exist only for light $Y$. Theorem~\ref{s3:thmT16s},
  with $t=t_Y=\lceil\lambda^{1.6}\rceil\le2\lambda^{1.6}$ and
  $\rho\ge1/(12L_Y^5)$, requires $s\ge2^{135}tL_Y^{28}\rho^{-5}$, and
  $2^{135}tL_Y^{28}\rho^{-5}\le2^{136}12^5\lambda^{1.6}L_Y^{53}$.
  With $s=s_Y/(8\klend)\ge\lambda^{100}/(16\bar k)$ and $L_Y\le2\lambda$,
  it suffices that
  $\lambda^{100}\ge2^{4+136+53}12^5\bar k\lambda^{54.6}$, i.e.\
  $\lambda^{45.4}\ge2^{193}\cdot12^5\cdot\bar k$. By row~9,
  $\bar k\le\lambda^{3.3}$, so the sufficient inequality of column~3,
  $\lambda^{42.1}\ge2^{193}\cdot12^5$, implies it. (Row~9 holds at
  $\mu$ under \Gc{1}; so column~3 implies column~2 under \Gc{1}.)
\item[5.] Theorem~\ref{s3:thmT16s} with $t=\tJS_l\le3M$ and
  $\rho=\rho_l\ge M^{-4}$ requires $s\ge2^{135}\cdot3M\cdot L_Y^{28}M^{20}$.
  With $s\ge\lambda^{100}/(16\bar k)$, $L_Y\le2\lambda$ and $M\le\bar M$,
  this becomes $\lambda^{100}\ge3\cdot2^{4+135+28}\bar k\bar M^{21}\lambda^{28}$.
\item[6.] The failure bound of Theorem~\ref{s3:thmT16s} is
  $2^{86}tL^{19}\rho^{-3}N^{-3}$ per class, with $N:=|V(Y)|$. Summed over
  the at most $\frac43M^2$ JS classes, with $t\le3M$, $\rho^{-3}\le M^{12}$
  and $L_Y\le2\lambda$, it is at most $2^{107}\lambda^{19}M^{15}N^{-3}$.
  This is at most $N^{-2}/4$ if $N\ge2^{109}\lambda^{19}M^{15}$, which by
  (B4) follows from $\lambda^{84}\ge2^{110}\bar M^{15}$. This line
  needs \Gc{1}: the class count and the multiplicity grow with $M$,
  which is bounded only by $\bar M=(A\mu)^{2A}$, not by a fixed power of
  $\lambda$ for moderate $\mu$. It
  does not follow from $C'>34$ alone (\RTb{I11}).
\item[7.] There are at most $12L_Y^3$ U-classes, each with
  $t=t_Y\le2\lambda^{1.6}$ and $\rho^{-3}\le1728L_Y^{15}$. The sum is at most
  $12L_Y^3\cdot2^{86}\cdot2\lambda^{1.6}L_Y^{19}\cdot12^3L_Y^{15}N^{-3}
  =12^4\cdot2^{87}\lambda^{1.6}L_Y^{37}N^{-3}\le12^4\cdot2^{124}\lambda^{38.6}N^{-3}$.
  This is at most $N^{-2}/4$ if $N\ge12^4\cdot2^{126}\lambda^{38.6}$,
  which by (B4) follows from
  $\lambda^{103}/2\ge12^4\cdot2^{126}\lambda^{38.6}$, i.e.\
  $\lambda^{64.4}\ge2^{127}\cdot12^4$.
\item[8.] By (B3), $s_r/4\ge\lambda^{100}/4$ and $s_r/2\ge\lambda^{100}/2$.
  Also $\kown\le L_Y\le2\lambda$, so $s'\ge\lambda^{99}/32$, and
  $\log L_Y\le\mu+1$. The right sides are increasing in $L_Y$.
  Substituting $L_Y\le2\lambda$ gives the four inequalities.
\item[9.--10.] By the first three inequalities of (i),
  $\klend(Y)\le\bar k$. So $\bar k\le\lambda^{3.3}$ suffices for row~9,
  and, since $p_Y=1/(2\klend)\ge1/(2\bar k)$, $2\bar k\le\lambda^4$
  suffices for row~10.
\item[11.] $\lambda^{95}/(8M_l^2)\ge2^{10}M_l^{10}$ is equivalent to
  $\lambda^{95}\ge2^{13}M_l^{12}$; use (B2).
\item[12.--13.] Row~12 is (B2) together with the inequality shown. Row~13
  is the condition \Gstar of Condition~\ref{s1:condGamma} itself.
\end{itemize}
\end{proof}

The next lemma shows that every column-3 inequality of
Table~\ref{s3:tabCOLJV} holds for all sufficiently large $\mu$. It has no
hypothesis. It is used only in the proof of Lemma~\ref{s7:lemGammaSat}, to show
that item~(f) of \Gc{1} can be satisfied.

\begin{lemma}[eventual form of the COL-JV table]\label{s3:lemCOLJVev}
Call an inequality in a real variable $\mu$ \emph{of type (E)} if it reads
$a\mu-b-c\log(A\mu)\ge0$ with reals $a>0$ and $b,c\ge0$ (here $A=105$ and
$\log=\log_2$).
\begin{itemize}
\item[(i)] An inequality of type (E) holds for every
  $\mu\ge\max\{1,v_0^2\}$, where
  $v_0\defeq\bigl(c+(c^2+a(b+c\log A))^{1/2}\bigr)/a$. In particular, it holds
  for all sufficiently large $\mu$.
\item[(ii)] Let $\mu\ge6$ be real, and put $\lambda\defeq2^\mu$,
  $\bar M\defeq(A\mu)^{2A}$ and $\bar k\defeq192\lambda^3+\frac43\bar M^2$. Let
  $1\le i\le13$, and let $(a_i,b_i,c_i)$ be the triple in column~4 of row~$i$
  of Table~\ref{s3:tabCOLJV}. If $a_i\mu-b_i-c_i\log(A\mu)\ge0$, then every
  inequality in column~3 of row~$i$ holds at $\lambda$.
\item[(iii)] Consequently, for each row of Table~\ref{s3:tabCOLJV}, every
  inequality in column~3 holds at $\lambda=2^\mu$ for all sufficiently large
  real $\mu$.
\end{itemize}
The lemma concerns explicit functions of the real variable $\mu$ only. It
involves no graph, no run of the hierarchy and no condition on $\Dstar$.
\deps{none (the inequalities are those printed in Table~\ref{s3:tabCOLJV},
which is not a statement).}
\end{lemma}

\begin{proof}
(i) Let $\mu\ge\max\{1,v_0^2\}$ and $v\defeq\mu^{1/2}$, so that $v\ge1$ and
$v\ge v_0$. For $v\ge1$ we have $2^v\ge1+v>v$: equality $2^v=1+v$ holds at
$v=1$, and for $v\ge1$ the derivative $2^v\ln2\ge2\ln2>1$ of the left side
exceeds the derivative $1$ of the right side. Hence $\log v<v$, and
$\log\mu=2\log v\le2v$. As $c\ge0$,
\[
a\mu-b-c\log(A\mu)\ =\ av^2-(b+c\log A)-c\log\mu\ \ge\ av^2-2cv-(b+c\log A).
\]
The right side is a quadratic polynomial in $v$ with leading coefficient
$a>0$. Its discriminant $4c^2+4a(b+c\log A)$ is non-negative, as $b,c\ge0$ and
$\log A>0$, and its larger root is $v_0$. So it is non-negative for every
$v\ge v_0$, and the inequality of type (E) holds at $\mu$.

(ii) We use three elementary bounds, valid for $\mu\ge6$. First,
$\log\mu\le\log(A\mu)$ and $\log(\mu+1)\le1+\log\mu\le1+\log(A\mu)$ (as
$\mu+1\le2\mu$), and $\log(A\mu)\ge0$. Second, $6\mu+20\le2^\mu$ (at $\mu=6$
this reads $56\le64$, and for $\mu\ge6$ the derivative $2^\mu\ln2\ge44$ of the
right side exceeds the derivative $6$ of the left side), so
$\log(2^\mu+6\mu+20)\le\log(2^{\mu+1})=\mu+1$; and
$(2\lambda)^6+1\le2\cdot(2\lambda)^6=2^7\lambda^6$. Third, both summands of
$\bar k$ are at least $1$ (as $\lambda\ge1$ and $A\mu\ge1$), and
$x+y\le2\max\{x,y\}\le2xy$ for $x,y\ge1$; so
$\bar k\le2\cdot192\lambda^3\cdot\frac43\bar M^2=2^9\lambda^3\bar M^2$ and
\[
\log\bar k\ \le\ 3\mu+4A\log(A\mu)+9 .
\]
Also $\log\bar M=2A\log(A\mu)$. For an inequality of column~3 put
$\varphi\defeq\log(\text{left side})-\log(\text{right side})$; the inequality
holds if and only if $\varphi\ge0$ (in row~1, where it is strict, if and only
if $\varphi>0$). The bounds above give, row by row:
\begin{itemize}
\item[1.] $\varphi=102.5\mu-\log257-102\log(2^\mu+6\mu+20)\ge102.5\mu-\log257-102(\mu+1)\ge0.5\mu-111$,
  as $\log257<9$; so $\varphi\ge(0.5\mu-112)+1>0$ if $0.5\mu-112\ge0$. The crude
  form $\lambda^{1/2}>2^{111}$ reads $\mu>222$, which also holds if
  $0.5\mu-112\ge0$, since then $\mu\ge224$;
\item[2.] $\varphi=103\mu-13\log\bar M=103\mu-26A\log(A\mu)$;
\item[3.] $\varphi=99\mu-\log640-\log\bar k\ge96\mu-19-4A\log(A\mu)$, as
  $\log640<10$;
\item[4.] $\varphi=42.1\mu-193-5\log12\ge42.1\mu-211$, as $5\log12<18$;
\item[5.] $\varphi=72\mu-\log3-167-\log\bar k-21\log\bar M\ge
  69\mu-178-46A\log(A\mu)$, as $\log3<2$;
\item[6.] $\varphi=84\mu-110-15\log\bar M=84\mu-110-30A\log(A\mu)$;
\item[7.] $\varphi=64.4\mu-127-4\log12\ge64.4\mu-142$, as $4\log12<15$;
\item[8.] put $E\defeq58\mu-194-\log(A\mu)$. The five inequalities of row~8
  hold if the following five differences are non-negative:
  $58\mu-194\ge E$;
  $62\mu-186-\log(\mu+1)\ge62\mu-187-\log(A\mu)=E+4\mu+7$;
  $62\mu-190-\log(\mu+1)\ge62\mu-191-\log(A\mu)=E+4\mu+3$;
  $58\mu-191\ge58\mu-194\ge E$; and
  $99\mu-6-\log((2\lambda)^6+1)\ge99\mu-6-(7+6\mu)=93\mu-13>0$. So all five
  hold if $E\ge0$;
\item[9.] $\varphi=3.3\mu-\log\bar k\ge0.3\mu-9-4A\log(A\mu)$;
\item[10.] $\varphi=4\mu-1-\log\bar k\ge\mu-10-4A\log(A\mu)$;
\item[11.] $\varphi=95\mu-13-12\log\bar M=95\mu-13-24A\log(A\mu)$;
\item[12.] $\varphi=1.6\mu-\log\bar M=1.6\mu-2A\log(A\mu)$;
\item[13.] $\varphi=36\mu-240-46A\log(A\mu)$.
\end{itemize}
In each row the lower bound obtained is $a_i\mu-b_i-c_i\log(A\mu)$ for the
triple $(a_i,b_i,c_i)$ of column~4 (in row~1 it is this quantity plus $1$, and
in row~8 it is at least $E=a_8\mu-b_8-c_8\log(A\mu)$ for the four inequalities
that depend on it). So if $a_i\mu-b_i-c_i\log(A\mu)\ge0$, every inequality in
column~3 of row~$i$ holds at $\lambda$. This proves (ii).

(iii) Fix a row~$i$, and let $v_0$ be as in (i) for
$(a,b,c)=(a_i,b_i,c_i)$. By (i), the inequality
$a_i\mu-b_i-c_i\log(A\mu)\ge0$ holds for every $\mu\ge\max\{1,v_0^2\}$; by (ii),
every inequality in column~3 of row~$i$ holds at every
$\mu\ge\max\{6,v_0^2\}$. This threshold is explicit, but its numerical value is
never computed: only its existence is used.
\end{proof}

\begin{lemma}[Lemma COL]\label{s3:lemCOL}
Assume condition~\ref{s1:condG1}. Let $Y$ be an ancestor of round $r$, with the stage-1
lending data of Definition~\ref{s3:defCOL}. If $r\ge R-1$, then
$\klend(Y)=0$ and every statement below about lent classes is vacuous.
\begin{itemize}
\item[(a)] $\Own_Y$ and $\Lend_Y$ are $(\varepsilon_Y,s_Y/4)$-expanders
  on $V(Y)$. Every lent class, as a graph on $V(Y)$, is an
  $(\varepsilon_Y,\,s_Y/(8\klend(Y)))$-expander. If $Y$ is light, every
  own class $\vR_{j,c}$, $\vM$, as a graph on $V(Y)$, is a
  $(2^{-6},\,s_r/(16\kown))$-expander.
\item[(b)] For all $r+2\le l\le R$ and $0\le j<\KJS_l$, the class
  $\LJS_{Y,l,j}$ is $(2^{12}L_Y^4,\tJS_l)$-path connected through
  $T_j(Y,l)$. By Lemma~\ref{s3:lemMonotone}(iii), one application routes
  any multiset of pairs of distinct vertices in which every vertex lies
  in at most $\tJS_l$ pairs. For instance, the union of the junction-$j$
  pair multisets of several systems, with multiplicities counted over
  the union, is joined by pairwise edge-disjoint paths of
  $\LJS_{Y,l,j}$ of length at most $2^{12}L_Y^4$ with interiors in
  $T_j(Y,l)$.
\item[(c)] Let $Y$ be light. Let $(V_{l,c,\uslot})_{(l,c,\uslot)\in\IU(Y)}$
  be random subsets of $V(Y)$, independent of the stage-1 lending data
  of $Y$, such that each $V_{l,c,\uslot}$ is a
  $\rho_{l,c,\uslot}$-random subset of $V(Y)$ with
  $\rho_{l,c,\uslot}\ge1/(12L_Y^5)$. The family may be dependent across
  indices. Then, with probability at least $1-|V(Y)|^{-2}/2$ over the
  lending data and the sets, every U-lent class $\LU_{Y,l,c,\uslot}$ is
  $(2^{12}L_Y^4,t_Y)$-path connected through $V_{l,c,\uslot}$, where
  $t_Y=\lceil\lambda_r^{1.6}\rceil$ (Definition~\ref{s3:defCOL}).
\item[(e)] The own-device thresholds hold. If $Y$ is standalone, then
  $s_Y/4\ge2^{150}L_Y^{42}$ and $s_Y/4\ge2^{146}L_Y^{38}\log L_Y$. On (a),
  $\Own_Y$ is therefore a spanning $(2^{-5},s_r/4)$-expander of $V(Y)$
  whose robustness meets the thresholds of the transparent P-vortex
  and of VX$^+$ (the second of these is not used: VX$^+$ is applied only to
  demoted light parts, Lemma~\ref{s5:lemDemoted}; cf.\
  Remark~\ref{s4:remTPVVX}). If $Y$ is light, then
  $s_r/2\ge2^{151}L_Y^{38}\log L_Y$; this is the threshold for VX$^+$ at
  expansion $2^{-6}$ on $H_Y=X_Y$, which is a
  $(2^{-6},s_r/2)$-expander. Moreover, on (a) every own class is a
  $(2^{-6},s')$-expander with $s':=s_r/(16\kown)\ge2^{145}L_Y^{41}$ and
  $2\lceil L_Y^6\rceil\le s'$.
\item[(g)] (COL(g).) The consumers listed in Definition~\ref{s3:defCOL}
  are exclusive. The own classes partition $\Own_Y$, and if $r\le R-2$,
  the lent classes are pairwise edge-disjoint and partition $\Lend_Y$. So
  no edge of $H_Y$ is available to two consumers, and every lent edge
  that its consumer does not use is junk for $Y$'s own device. If
  $r\ge R-1$, then $\Lend_Y$ has no classes, and all of $\Lend_Y$ is junk
  for $Y$'s own device (Definition~\ref{s3:defCOL}, Consumers).
\end{itemize}
Items (a), (b) and (e) hold simultaneously with probability at least
$1-|V(Y)|^{-2}/2$ over the stage-1 lending data of $Y$. For every family
as in (c), items (a), (b), (c) and (e) hold simultaneously with
probability at least $1-|V(Y)|^{-2}$. Item (g) always holds. The failure
unions are bounded through rows~6 and~7 of Table~\ref{s3:tabCOLJV}, which
need \Gc{1} (\RTb{I11}). There are no items (d) and (f); the letters
follow the references used in later sections.
\deps{\ref{s3:defCOL}, \ref{s3:lemL15p}, \ref{s3:thmT16s}, \ref{s3:lemMonotone}, \ref{s3:lemCOLJV}, \ref{s2:lemTower}, \ref{s2:propStructure}, \ref{s2:defAncestors}, \ref{s1:condG1}}
\fixes{\RTb{I11}, \textsf{CR1-PV} (item (c) with $t_Y$)}
\end{lemma}

\begin{proof}
Write $N:=|V(Y)|$, $L:=L_Y$, $k:=\klend(Y)$, $\lambda:=\lambda_r$ and
$\mu:=\log\lambda$. By Proposition~\ref{s2:propStructure}(i),
$H_Y$ is an $(\varepsilon_Y,s_Y)$-expander on $V(Y)$. Here
$(\varepsilon_Y,s_Y)=(2^{-5},s_r)$ if $Y$ is standalone and
$(2^{-6},s_r/2)$ if $Y$ is light
(Definition~\ref{s2:defAncestors}), and $s_r\ge\lambda^{100}$
(Lemma~\ref{s2:lemTower}(b)). The standing bounds (B1) and (B3)--(B6) of
the proof of Lemma~\ref{s3:lemCOLJV} are available in every round, and
(B2), (B7) when $r\le R-2$. Rows~1, 3, 8 and~9 of
Table~\ref{s3:tabCOLJV} hold for every $r\le R$, and the other rows
involving $Y$ hold for $r\le R-2$ (Lemma~\ref{s3:lemCOLJV}(ii)); for
$r\ge R-1$ only rows~3, 8 and~9 are used below. In particular
$N\ge\lambda^{103}/2$ and $L\le2\lambda$.

\emph{(a)} Stage (i) is a uniform $2$-colouring of $E(H_Y)$. By
Lemma~\ref{s3:lemL15p} with $k=2$, which needs $s_Y\ge80L$ (row~3 of
Table~\ref{s3:tabCOLJV}), $\Own_Y$ and $\Lend_Y$ are
$(\varepsilon_Y,s_Y/4)$-expanders on $V(Y)$, except with probability at
most $4N^{-5}$.

Condition on stage (i) with $\Lend_Y$ such an expander. Stage (ii) is
then a uniform $k$-colouring of $\Lend_Y$. By Lemma~\ref{s3:lemL15p},
which needs $s_Y/4\ge40kL$ (row~3), every lent class is an
$(\varepsilon_Y,s_Y/(8k))$-expander on $V(Y)$, except with probability
at most $2kN^{-5}$.

If $Y$ is light, condition on $\Own_Y$ being a
$(2^{-6},s_r/8)$-expander. Stage (iii) is a uniform $\kown$-colouring.
Lemma~\ref{s3:lemL15p}, which needs $s_r/8\ge40\kown L$ (row~3), makes
every own class an expander with parameters $2^{-6}$ and
$s_r/(16\kown)$, except with probability at most $2\kown N^{-5}$.

Using row~9 ($k\le\lambda^{3.3}$), $\kown\le L\le2\lambda$ and
$N^3\ge\lambda^{309}/8$, the total failure probability of (a) is at most
\[
2(2+k+\kown)N^{-5}\le2(2+\lambda^{3.3}+2\lambda)N^{-5}\le N^{-2}/4 .
\]

\emph{(b)} Condition on a stage-(i)/(ii) outcome in which (a) holds. The
labels are independent of the colouring, so for fixed $(l,j)$ the set
$T_j(Y,l)$ is still a $\rho_l$-random subset of $V(Y)$. The class
$\LJS_{Y,l,j}$ is now a fixed $N$-vertex $(\varepsilon_Y,s)$-expander
with $s:=s_Y/(8k)$ and $\varepsilon_Y\in[2^{-7},1]$. Apply
Theorem~\ref{s3:thmT16s} with $t:=\tJS_l$ and $\rho:=\rho_l$. Its
hypotheses hold:
\begin{itemize}
\item $s\ge2^{135}tL^{28}\rho^{-5}$ by row~5;
\item $\rho N\ge L^2$, since by row~2 at $\lambda=\lambda_r$
  (Lemma~\ref{s3:lemCOLJV}(ii)), $\bar M^4\le\lambda^{31.7}$, so
  $\rho_lN\ge\bar M^{-4}\lambda^{103}/2\ge\lambda^{71}/2\ge4\lambda^2\ge L^2$.
\end{itemize}
Each class fails with probability at most
$2^{86}\tJS_lL^{19}\rho_l^{-3}N^{-3}$. By row~6, the sum over $\IJS(Y)$
is at most $N^{-2}/4$. The final sentence of (b) is
Lemma~\ref{s3:lemMonotone}(iii).

\emph{(c)} Condition as in (b). The sets $V_{l,c,\uslot}$ are independent
of the colouring. For a fixed index, apply Theorem~\ref{s3:thmT16s} to
the fixed class $\LU_{Y,l,c,\uslot}$, an $N$-vertex
$(2^{-6},s_r/(16k))$-expander, with $t:=t_Y=\lceil\lambda^{1.6}\rceil$
(so $1\le t\le2\lambda^{1.6}$) and
$\rho:=\rho_{l,c,\uslot}\ge1/(12L^5)$. Its hypotheses hold:
\begin{itemize}
\item $s_r/(16k)\ge2^{135}tL^{28}\rho^{-5}$ by row~4 (with row~9), since
  $\rho^{-5}\le(12L^5)^5$;
\item $\rho N\ge N/(12L^5)\ge\lambda^{103}/(24(2\lambda)^5)\ge L^2$.
\end{itemize}
Each index fails with probability at most
$2^{86}tL^{19}(12L^5)^3N^{-3}$. The theorem is applied to one index at
a time, so dependence across indices is irrelevant. By row~7, the union
over $\IU(Y)$ is at most $N^{-2}/4$. Together with the failure
probability of (a), (c) fails with probability at most $N^{-2}/2$.

\emph{(e)} The inequalities are row~8 of Table~\ref{s3:tabCOLJV}, with
$s_Y=s_r$ for standalone $Y$. On (a), $\Own_Y$ is a
$(2^{-5},s_r/4)$-expander if $Y$ is standalone, and the own classes are
$(2^{-6},s_r/(16\kown))$-expanders if $Y$ is light. Also $H_Y=X_Y$ is a
$(2^{-6},s_r/2)$-expander for light $Y$, by
Proposition~\ref{s2:propStructure}(i).

\emph{(g)} This holds by construction. If $r\le R-2$, each edge of
$\Lend_Y$ receives exactly one index; if $r\ge R-1$, $\Lend_Y$ has no
classes and is junk for $Y$'s own device. Each edge of $\Own_Y$ receives
exactly one own label.

\emph{Probabilities.} By the bounds above, (a) or (b) fails with
probability at most $N^{-2}/4+N^{-2}/4=N^{-2}/2$, and (e) holds on (a).
Adding the U-union of (c) gives at most $3N^{-2}/4\le N^{-2}$.
\end{proof}
